% Chapter 8. All quoted numbers come from chapters/ch08-numbers.tex, written by code/ch08/run_all.py.
% Generated by code/ch08/run_all.py. Do not edit by hand.
\newcommand{\acAbs}{0.26}
\newcommand{\acAbsAbove}{100}
\newcommand{\acAbsHundred}{0.10}
\newcommand{\acLags}{100}
\newcommand{\acRet}{0.01}
\newcommand{\acRetOut}{18}
\newcommand{\acSq}{0.23}
\newcommand{\acSqAbove}{100}
\newcommand{\acSqHundred}{0.028}
\newcommand{\caCritEmp}{0.087}
\newcommand{\caCritNor}{0.085}
\newcommand{\caCritT}{0.087}
\newcommand{\caEwEs}{0.55}
\newcommand{\caEwPc}{\ensuremath{<}0.001}
\newcommand{\caEwPk}{\ensuremath{<}0.001}
\newcommand{\caEwRate}{1.79}
\newcommand{\caEwRed}{2}
\newcommand{\caEwX}{266}
\newcommand{\caEwXs}{500}
\newcommand{\caExpected}{149}
\newcommand{\caExpectedEs}{372}
\newcommand{\caFhsEs}{0.17}
\newcommand{\caFhsPc}{0.068}
\newcommand{\caFhsPk}{\ensuremath{<}0.001}
\newcommand{\caFhsRate}{1.32}
\newcommand{\caFhsRed}{0}
\newcommand{\caFhsX}{197}
\newcommand{\caFhsXs}{416}
\newcommand{\caFirst}{20 June 1967}
\newcommand{\caGnEs}{0.68}
\newcommand{\caGnPc}{0.002}
\newcommand{\caGnPk}{\ensuremath{<}0.001}
\newcommand{\caGnRate}{2.01}
\newcommand{\caGnRed}{4}
\newcommand{\caGnX}{300}
\newcommand{\caGnXs}{545}
\newcommand{\caGtEs}{0.46}
\newcommand{\caGtPc}{0.010}
\newcommand{\caGtPk}{\ensuremath{<}0.001}
\newcommand{\caGtRate}{1.51}
\newcommand{\caGtRed}{0}
\newcommand{\caGtX}{225}
\newcommand{\caGtXs}{516}
\newcommand{\caHsEs}{0.34}
\newcommand{\caHsPc}{\ensuremath{<}0.001}
\newcommand{\caHsPk}{\ensuremath{<}0.001}
\newcommand{\caHsRate}{1.47}
\newcommand{\caHsRed}{4}
\newcommand{\caHsX}{219}
\newcommand{\caHsXs}{460}
\newcommand{\caHslEs}{0.30}
\newcommand{\caHslPc}{\ensuremath{<}0.001}
\newcommand{\caHslPk}{\ensuremath{<}0.001}
\newcommand{\caHslRate}{1.46}
\newcommand{\caHslRed}{10}
\newcommand{\caHslX}{218}
\newcommand{\caHslXs}{454}
\newcommand{\caMaxFhs}{9}
\newcommand{\caMaxGt}{10}
\newcommand{\caMaxHs}{15}
\newcommand{\caMaxNor}{21}
\newcommand{\caN}{14,897}
\newcommand{\caNorEs}{0.67}
\newcommand{\caNorPc}{\ensuremath{<}0.001}
\newcommand{\caNorPk}{\ensuremath{<}0.001}
\newcommand{\caNorRate}{2.16}
\newcommand{\caNorRed}{16}
\newcommand{\caNorX}{322}
\newcommand{\caNorXs}{512}
\newcommand{\caNuLast}{6.7}
\newcommand{\caNullMeanEmp}{\ensuremath{-}0.001}
\newcommand{\caNullMeanNor}{0.000}
\newcommand{\caNullMeanT}{0.000}
\newcommand{\caNullSdEmp}{0.054}
\newcommand{\caNullSdNor}{0.052}
\newcommand{\caNullSdT}{0.053}
\newcommand{\caOctFhs}{3.6}
\newcommand{\caOctGt}{3.7}
\newcommand{\caOctHs}{7.2}
\newcommand{\caTableRows}{Normal, 250 days & \caNorX & \caNorPk & \caNorPc & \caNorRed & \caNorEs & \cyNorX & \cyNorPk & \cyNorPc \\
Historical, 250 days & \caHsX & \caHsPk & \caHsPc & \caHsRed & \caHsEs & \cyHsX & \cyHsPk & \cyHsPc \\
Historical, 1,000 days & \caHslX & \caHslPk & \caHslPc & \caHslRed & \caHslEs & \cyHslX & \cyHslPk & \cyHslPc \\
EWMA, normal & \caEwX & \caEwPk & \caEwPc & \caEwRed & \caEwEs & \cyEwX & \cyEwPk & \cyEwPc \\
GARCH, normal & \caGnX & \caGnPk & \caGnPc & \caGnRed & \caGnEs & \cyGnX & \cyGnPk & \cyGnPc \\
GARCH, $t$ & \caGtX & \caGtPk & \caGtPc & \caGtRed & \caGtEs & \cyGtX & \cyGtPk & \cyGtPc \\
Filtered historical & \caFhsX & \caFhsPk & \caFhsPc & \caFhsRed & \caFhsEs & \cyFhsX & \cyFhsPk & \cyFhsPc \\}
\newcommand{\cfDownHi}{74}
\newcommand{\cfHsThree}{5.8}
\newcommand{\cfLossNine}{16.6}
\newcommand{\cfNegHi}{88}
\newcommand{\cfZnine}{\ensuremath{-}16.6}
\newcommand{\crAutDays}{63}
\newcommand{\crAutFhs}{3}
\newcommand{\crAutHs}{9}
\newcommand{\crDays}{756}
\newcommand{\crDropDate}{29 September 2008}
\newcommand{\crDropLoss}{8.25}
\newcommand{\crExpected}{7.6}
\newcommand{\crFhs}{14}
\newcommand{\crFhsPeak}{16 October 2008}
\newcommand{\crHs}{22}
\newcommand{\crHsDrop}{28 September 2009}
\newcommand{\crHsMax}{8.25}
\newcommand{\crHsPeak}{2 December 2008}
\newcommand{\cyEwPc}{0.207}
\newcommand{\cyEwPk}{0.012}
\newcommand{\cyEwX}{55}
\newcommand{\cyExpected}{39}
\newcommand{\cyFhsPc}{0.552}
\newcommand{\cyFhsPk}{0.031}
\newcommand{\cyFhsX}{26}
\newcommand{\cyFirst}{27 April 2011}
\newcommand{\cyGnPc}{0.492}
\newcommand{\cyGnPk}{0.151}
\newcommand{\cyGnX}{30}
\newcommand{\cyGtPc}{0.664}
\newcommand{\cyGtPk}{\ensuremath{<}0.001}
\newcommand{\cyGtX}{19}
\newcommand{\cyHsPc}{0.114}
\newcommand{\cyHsPk}{0.307}
\newcommand{\cyHsX}{45}
\newcommand{\cyHslPc}{0.155}
\newcommand{\cyHslPk}{0.019}
\newcommand{\cyHslX}{25}
\newcommand{\cyLast}{29 September 2026}
\newcommand{\cyN}{3,852}
\newcommand{\cyNorPc}{0.057}
\newcommand{\cyNorPk}{0.008}
\newcommand{\cyNorX}{56}
\newcommand{\ddAllYears}{4.4}
\newcommand{\ddCagr}{10.9}
\newcommand{\ddCalmar}{0.20}
\newcommand{\ddDownYears}{1.4}
\newcommand{\ddEmpMean}{9.5}
\newcommand{\ddLawMean}{12.8}
\newcommand{\ddLawTwenty}{17.5}
\newcommand{\ddLogMu}{10.33}
\newcommand{\ddLogVol}{16.26}
\newcommand{\ddMax}{55}
\newcommand{\ddPeak}{9 October 2007}
\newcommand{\ddRecover}{28 February 2012}
\newcommand{\ddSecond}{49}
\newcommand{\ddSecondPeak}{March 2000}
\newcommand{\ddSecondTrough}{October 2002}
\newcommand{\ddShareTen}{29}
\newcommand{\ddShareTwenty}{14}
\newcommand{\ddShareUnder}{89}
\newcommand{\ddSimFifty}{2}
\newcommand{\ddSimM}{4,000}
\newcommand{\ddSimMax}{67}
\newcommand{\ddSimMedian}{29}
\newcommand{\ddSimNinety}{41}
\newcommand{\ddSimThirty}{43}
\newcommand{\ddSimpleHi}{42}
\newcommand{\ddSimpleLo}{33}
\newcommand{\ddThird}{48}
\newcommand{\ddThirdPeak}{January 1973}
\newcommand{\ddThirdTrough}{October 1974}
\newcommand{\ddTrough}{9 March 2009}
\newcommand{\ddWinAboveNinety}{26}
\newcommand{\ddWinMax}{55}
\newcommand{\ddWinMedian}{37}
\newcommand{\ddWinMin}{20}
\newcommand{\ddWinN}{54}
\newcommand{\esGapN}{0.0002}
\newcommand{\esGapS}{0.0015}
\newcommand{\esHsN}{4.10}
\newcommand{\esHsS}{3.04}
\newcommand{\esKN}{159}
\newcommand{\esKS}{398}
\newcommand{\esNorN}{2.67}
\newcommand{\esNorS}{2.34}
\newcommand{\esRatioHist}{1.09}
\newcommand{\esRatioNormal}{1.005}
\newcommand{\esTailN}{158.970}
\newcommand{\esTailS}{397.425}
\newcommand{\esTtN}{4.38}
\newcommand{\esTtS}{3.12}
\newcommand{\evAlphaHi}{6.1}
\newcommand{\evAlphaLo}{2.9}
\newcommand{\evAlphaTail}{3.96}
\newcommand{\evBeta}{0.645}
\newcommand{\evEsF}{15.36}
\newcommand{\evEsN}{4.10}
\newcommand{\evEsT}{8.14}
\newcommand{\evHillEight}{2.7}
\newcommand{\evHillTwoHundred}{3.0}
\newcommand{\evHsVarF}{12.00}
\newcommand{\evHsVarN}{2.79}
\newcommand{\evHsVarT}{5.91}
\newcommand{\evNorVarF}{3.75}
\newcommand{\evNorVarN}{2.33}
\newcommand{\evNorVarT}{3.11}
\newcommand{\evNormalYears}{\ensuremath{5\times 10^{62}}}
\newcommand{\evNu}{791}
\newcommand{\evNuTop}{159}
\newcommand{\evSimBias}{0.4}
\newcommand{\evSimHsSd}{6.3}
\newcommand{\evSimM}{400}
\newcommand{\evSimSd}{5.3}
\newcommand{\evSimXi}{0.194}
\newcommand{\evSimXiSd}{0.044}
\newcommand{\evU}{1.53}
\newcommand{\evVarF}{11.23}
\newcommand{\evVarN}{2.81}
\newcommand{\evVarT}{5.83}
\newcommand{\evXi}{0.2525}
\newcommand{\evXiMax}{0.34}
\newcommand{\evXiMaxQ}{98}
\newcommand{\evXiMin}{0.16}
\newcommand{\evXiSe}{0.0445}
\newcommand{\evXiSeTop}{0.10}
\newcommand{\evYears}{200}
\newcommand{\ewCorr}{0.975}
\newcommand{\ewLLdiff}{214}
\newcommand{\fcKaDate}{20 October 1987}
\newcommand{\fcKaNow}{93.4}
\newcommand{\fcKaReal}{22}
\newcommand{\fcKaYear}{58}
\newcommand{\fcKbDate}{20 November 2008}
\newcommand{\fcKbNow}{65.9}
\newcommand{\fcKbReal}{33}
\newcommand{\fcKbYear}{42}
\newcommand{\fcKcDate}{16 March 2020}
\newcommand{\fcKcNow}{83.4}
\newcommand{\fcKcReal}{29}
\newcommand{\fcKcYear}{53}
\newcommand{\fcKdDate}{14 July 2017}
\newcommand{\fcKdNow}{9.3}
\newcommand{\fcKdReal}{12}
\newcommand{\fcKdYear}{15}
\newcommand{\gaAlphaN}{0.101}
\newcommand{\gaAlphaT}{0.095}
\newcommand{\gaBetaN}{0.889}
\newcommand{\gaBetaT}{0.902}
\newcommand{\gaFourth}{1.001}
\newcommand{\gaHalfN}{70}
\newcommand{\gaHalfT}{210}
\newcommand{\gaLR}{757}
\newcommand{\gaLogLikN}{\ensuremath{-}19,553.5}
\newcommand{\gaLogLikT}{\ensuremath{-}19,175.1}
\newcommand{\gaLrVolN}{17.5}
\newcommand{\gaLrVolT}{24.1}
\newcommand{\gaNQThou}{\ensuremath{-}3.09}
\newcommand{\gaNu}{6.7}
\newcommand{\gaNuFine}{6.665}
\newcommand{\gaOmegaN}{0.0120}
\newcommand{\gaOmegaT}{0.0076}
\newcommand{\gaPersN}{0.9901}
\newcommand{\gaPersT}{0.9967}
\newcommand{\gaResAcSq}{0.01}
\newcommand{\gaResKurt}{2.5}
\newcommand{\gaResQ}{2.70}
\newcommand{\gaResQThou}{\ensuremath{-}4.61}
\newcommand{\gaResSkew}{\ensuremath{-}0.49}
\newcommand{\gaTQ}{\ensuremath{-}2.54}
\newcommand{\gaTQRaw}{\ensuremath{-}3.04}
\newcommand{\gaTQThou}{\ensuremath{-}4.11}
\newcommand{\kpHi}{6}
\newcommand{\kpLo}{1}
\newcommand{\kpSixLR}{3.56}
\newcommand{\kpSixP}{0.059}
\newcommand{\prAltKurt}{6}
\newcommand{\prAltSkew}{\ensuremath{-}1}
\newcommand{\prEta}{3.38}
\newcommand{\prFirst}{July 1963}
\newcommand{\prHtAlpha}{0.28}
\newcommand{\prHtBeta}{1.24}
\newcommand{\prHtIR}{0.03}
\newcommand{\prHtMsq}{6.45}
\newcommand{\prHtOmega}{1.37}
\newcommand{\prHtSR}{0.42}
\newcommand{\prHtSortino}{0.63}
\newcommand{\prHtTE}{11.0}
\newcommand{\prHtTreynor}{7.45}
\newcommand{\prHtVol}{22.0}
\newcommand{\prKurt}{1.7}
\newcommand{\prLast}{August 2026}
\newcommand{\prMktMean}{7.22}
\newcommand{\prMktOmega}{1.42}
\newcommand{\prMktSortino}{0.69}
\newcommand{\prMktTreynor}{7.22}
\newcommand{\prMktVol}{15.4}
\newcommand{\prSR}{0.47}
\newcommand{\prSRlo}{0.46}
\newcommand{\prSRm}{0.135}
\newcommand{\prSeIid}{0.126}
\newcommand{\prSeNaive}{0.126}
\newcommand{\prSeNn}{0.131}
\newcommand{\prSeTen}{0.33}
\newcommand{\prSkew}{\ensuremath{-}0.49}
\newcommand{\prT}{758}
\newcommand{\prYearsOne}{4.9}
\newcommand{\prYearsOneAlt}{5.8}
\newcommand{\prYearsTwo}{20}
\newcommand{\pwFiveFifteen}{42}
\newcommand{\pwOneFifteen}{11}
\newcommand{\pwOneTwo}{24}
\newcommand{\pwSizeOne}{9.5}
\newcommand{\pwTenFifteen}{62}
\newcommand{\pwYearsFifteen}{16}
\newcommand{\pwYearsTwo}{5}
\newcommand{\saBandHi}{1}
\newcommand{\saBandLo}{0.5013}
\newcommand{\saBoth}{1.79}
\newcommand{\saEsOne}{90}
\newcommand{\saEsTwo}{100.8}
\newcommand{\saP}{0.9}
\newcommand{\smAcS}{0.61}
\newcommand{\smAcT}{0.04}
\newcommand{\smAcU}{0.03}
\newcommand{\smAnnVolS}{15.5}
\newcommand{\smAnnVolT}{17.2}
\newcommand{\smBetaSum}{0.87}
\newcommand{\smBetaZero}{0.40}
\newcommand{\smCorrU}{0.9999}
\newcommand{\smExMeanS}{7.18}
\newcommand{\smExMeanT}{7.22}
\newcommand{\smExMeanU}{7.24}
\newcommand{\smExVolS}{7.85}
\newcommand{\smExVolT}{15.44}
\newcommand{\smExVolU}{15.83}
\newcommand{\smMddS}{42}
\newcommand{\smMddT}{50}
\newcommand{\smMddU}{51}
\newcommand{\smPhi}{0.6}
\newcommand{\smPhiHat}{0.610}
\newcommand{\smRatio}{0.50}
\newcommand{\smSRS}{0.91}
\newcommand{\smSRT}{0.47}
\newcommand{\smSRU}{0.46}
\newcommand{\smSRloS}{0.49}
\newcommand{\smSRloT}{0.46}
\newcommand{\smSRloU}{0.45}
\newcommand{\smVarS}{6.3}
\newcommand{\smVarT}{11.0}
\newcommand{\smVarU}{11.1}
\newcommand{\smVolRatioTheory}{0.50}
\newcommand{\smVolS}{7.8}
\newcommand{\smVolT}{15.4}
\newcommand{\smVolU}{15.8}
\newcommand{\srOne}{2.79}
\newcommand{\srOneEs}{4.10}
\newcommand{\srScaled}{8.82}
\newcommand{\srScaledEs}{12.96}
\newcommand{\srTen}{8.57}
\newcommand{\srTenEs}{12.42}
\newcommand{\tdBest}{11.4}
\newcommand{\tdBestDate}{13 October 2008}
\newcommand{\tdBeyondFive}{44}
\newcommand{\tdBeyondFiveNormal}{0.009}
\newcommand{\tdBeyondFour}{93}
\newcommand{\tdBeyondFourNormal}{1.0}
\newcommand{\tdFirst}{1 July 1963}
\newcommand{\tdKurt}{15.6}
\newcommand{\tdLast}{31 August 2026}
\newcommand{\tdMeanAnn}{11.7}
\newcommand{\tdN}{15,897}
\newcommand{\tdNu}{3.0}
\newcommand{\tdSd}{1.02}
\newcommand{\tdSkew}{\ensuremath{-}0.49}
\newcommand{\tdVolAnn}{16.2}
\newcommand{\tdWorst}{17.4}
\newcommand{\tdWorstDate}{19 October 1987}
\newcommand{\tdWorstSd}{17}
\newcommand{\thFloor}{50}
\newcommand{\thLeftEs}{17.9}
\newcommand{\thLeftMp}{2.75}
\newcommand{\thLeftMpSample}{2.77}
\newcommand{\thLeftOm}{1.39}
\newcommand{\thLeftSR}{0.43}
\newcommand{\thLeftSo}{0.54}
\newcommand{\thLeftTail}{4.5\times10^{-12}}
\newcommand{\thLeftVar}{13.9}
\newcommand{\thMu}{6}
\newcommand{\thNormalEs}{10.2}
\newcommand{\thNormalMp}{3.12}
\newcommand{\thNormalMpSample}{3.10}
\newcommand{\thNormalOm}{1.37}
\newcommand{\thNormalSR}{0.43}
\newcommand{\thNormalSo}{0.68}
\newcommand{\thNormalVar}{8.8}
\newcommand{\thRightEs}{3.5}
\newcommand{\thRightMp}{3.39}
\newcommand{\thRightMpSample}{3.38}
\newcommand{\thRightOm}{1.43}
\newcommand{\thRightSR}{0.43}
\newcommand{\thRightSo}{1.01}
\newcommand{\thRightVar}{3.5}
\newcommand{\thSR}{0.43}
\newcommand{\thVol}{13.9}
\newcommand{\volMax}{99}
\newcommand{\volMaxDate}{17 March 2020}
\newcommand{\volMedian}{12.3}
\newcommand{\volMin}{5.8}
\newcommand{\vrCfN}{6.32}
\newcommand{\vrCfS}{3.25}
\newcommand{\vrHsN}{2.79}
\newcommand{\vrHsS}{2.03}
\newcommand{\vrHsThree}{5.91}
\newcommand{\vrNorN}{2.33}
\newcommand{\vrNorS}{1.95}
\newcommand{\vrNorThree}{3.11}
\newcommand{\vrTtN}{2.80}
\newcommand{\vrTtS}{1.93}
\newcommand{\xaEs}{0.281}
\newcommand{\xaVar}{0.279}
\newcommand{\xaVarTen}{0.883}
\newcommand{\xbBoth}{0.16}
\newcommand{\xbEsOne}{80}
\newcommand{\xbEsTwo}{103.2}
\newcommand{\xbJoint}{7.84}
\newcommand{\xcHalf}{34.3}
\newcommand{\xcLr}{1.0}
\newcommand{\xcLrVol}{15.9}
\newcommand{\xcTen}{3.50}
\newcommand{\xcTenDay}{6.12}
\newcommand{\xcTenSqrt}{6.32}
\newcommand{\xdLR}{7.73}
\newcommand{\xdP}{0.0054}
\newcommand{\xdTail}{0.40}
\newcommand{\xeEs}{9.13}
\newcommand{\xeVar}{6.65}
\newcommand{\xfAnn}{0.52}
\newcommand{\xfAnnSe}{0.34}
\newcommand{\xfIid}{0.0918}
\newcommand{\xfNn}{0.0991}
\newcommand{\xfT}{1.51}
\newcommand{\xgInfl}{1.73}
\newcommand{\xgVol}{9.24}
\newcommand{\xhAw}{0.50}
\newcommand{\xhDisp}{2.45}
\newcommand{\xhDispAw}{2.01}
\newcommand{\xhEw}{2.00}
\newcommand{\xiLogThirty}{0.3567}
\newcommand{\xiMean}{18.75}
\newcommand{\xiThirty}{14.9}
\newcommand{\xjMsq}{6.0}
\newcommand{\xjTreynor}{7.27}
\newcommand{\znExpected}{2.5}
\newcommand{\znGreen}{89.2}
\newcommand{\znGreenThree}{13}
\newcommand{\znGreenTwo}{44}
\newcommand{\znRed}{0.025}
\newcommand{\znRedThree}{22}
\newcommand{\znRedTwo}{3}
\newcommand{\znYellow}{10.8}
\newcommand{\ztCrit}{0.0877}
\newcommand{\ztCritAn}{0.0887}
\newcommand{\ztCritSe}{0.0006}
\newcommand{\ztCritSeAn}{0.0008}
\newcommand{\ztM}{20,000}
\newcommand{\ztMean}{\ensuremath{-}0.0003}
\newcommand{\ztMeanSe}{0.0005}
\newcommand{\ztSd}{0.0532}
\newcommand{\ztSdAn}{0.0539}
\newcommand{\ztStreams}{10}
\newcommand{\ztT}{14,897}

\raggedbottom % this chapter only: short pages stay short; \flushbottom is restored at the end
\begin{chapterrefs}
\setcounter{chapter}{7}
\chapter{Risk and Performance Measurement}
\label{ch:risk}

A risk report for an endowment carries a handful of numbers: a volatility, a value at risk, perhaps
an expected shortfall, a worst drawdown, and a Sharpe ratio for each manager. Every one of them is an
estimate from a short and unrepresentative record, and every one rests on assumptions about how
returns behave from one day to the next that the data contradict. This chapter asks what each number
means, how it should be computed, and how much a record can say about whether it was right. It ends
with a test over almost sixty years of daily returns of the US stock market: of seven common ways to
set a one day value at risk, which ones produced the number of losses they promised?

The first half is about the size of losses. A risk measure should reward diversification, and value
at risk does not always do so, while expected shortfall does. Daily returns have tails far heavier
than the normal law and volatility that comes in clusters, so a measure computed from the whole
sample is too low in a crisis and too high in calm years. A GARCH\index{GARCH} model tracks the clusters, and
extreme value theory describes the far tail with a law that has only two parameters. Backtests count
the days on which the loss exceeded the forecast, and a little probability shows how little a year of
such counts can reveal. The second half is about returns per unit of risk. The Sharpe ratio and its
relatives rank managers, but they can be raised by selling insurance or by reporting smoothed values,
and their standard errors are large. Drawdowns measure what an investor actually lives through.
The last section before the case study turns to what a performance record must contain before any of
these numbers can be trusted.

\begin{keybox}[How to read this chapter]
The numbers in the text are computed by nine scripts in \texttt{code/ch08/} (\texttt{tails},
\texttt{garch}, \texttt{evt}, \texttt{backtest\_theory}, \texttt{ratios}, \texttt{drawdown},
\texttt{smoothing}, \texttt{case\_backtest} and \texttt{exercises}), and the tests in
\texttt{code/ch08/test\_*.py} check them, among other things against the GARCH estimates of the
\texttt{arch} package \parencite{arch}. An independent R script recomputes the moments, the
historical, normal and Cornish-Fisher measures, the GARCH likelihood with its own optimiser, the
generalised Pareto\index{generalised Pareto} fit, the backtest statistics of the case study, the ratios, the drawdowns and the
smoothing results. GNU Octave fits the GARCH model again with \texttt{fminsearch} and checks the
zone probabilities, the drawdown and the smoothing. The five listings printed in the chapter are
in \texttt{code/ch08/snippets/}; each prints figures quoted in the text, and
\texttt{code/ch08/test\_assets.py} checks them.

The data are the daily return of the US stock market and the one month Treasury bill rate from the
Kenneth R. French Data Library \parencite{frenchlib}, from \tdFirst{} to \tdLast, \tdN{} trading days;
the market return is the published excess return plus the bill rate. The library gives returns in
percent to two decimals, so daily returns are rounded to a hundredth of a percent. Monthly returns of
the market and of the ten industries of Chapter~7 come from the same library, and the
second series of the case study is the daily USD/CAD rate of the Bank of Canada used in Chapter~5.
A loss is minus a return. Value at risk and expected shortfall are quoted as positive losses, in
percent of the position. An annualised mean is 252 times a daily mean and an annualised volatility
$\sqrt{252}$ times a daily standard deviation, which treats daily returns as uncorrelated.

Random draws come from NumPy's default generator with base seed \texttt{8317} and a stream number
fixed in each script, so every run is repeatable.

Notation. $L$ is a loss with distribution function $F_L$, $\alpha$ a confidence level such as 0.99,
$\mathrm{VaR}_\alpha$ and $\mathrm{ES}_\alpha$ the value at risk and expected shortfall at that
level. $r_t$ is a daily return, $\sigma_t^2$ its variance given the past, and $z_t=(r_t-\mu)/\sigma_t$
the standardised residual. $\Phi$ and $\varphi$ are the standard normal distribution and density, and
$z_\alpha=\Phi^{-1}(\alpha)$. $\theta$ is a Sharpe ratio, $\beta$ a beta on the market, $T$ a number of
observations and $I_t$ the indicator of an exception, a day on which the loss exceeded the VaR.
A few letters are reused locally and redefined where they appear: $\beta$ is also the scale of the
generalised Pareto law, $s$ the scale of a $t$ law, $\rho$ a risk measure, an autocorrelation or a
risk aversion, and $\lambda$ a scale factor, the EWMA\index{EWMA} decay or a drawdown rate. The threshold of the
Omega ratio is $\tau$, and $\gamma_3$ and $\gamma_4$ are the skewness and excess kurtosis of returns
throughout.
\end{keybox}

\FloatBarrier
% =====================================================================================
\section{What should a risk measure do?}
\label{sec8:measures}

Volatility is the oldest measure of risk and still the most used. It treats a gain and a loss of the
same size alike, and it says nothing about the tail unless the law of returns is known. Downside
measures look only at losses. The \emph{value at risk}\index{value at risk} at level $\alpha$ is the loss that is exceeded
with probability at most $1-\alpha$,
\begin{equation}\label{eq8:var}
  \mathrm{VaR}_\alpha(L) = \inf\{\ell : \mathbb P(L>\ell)\le 1-\alpha\},
\end{equation}
so a one day 99\% VaR of 2\% says that on one day in a hundred the loss should exceed 2\% of the
position. It says nothing about how large the loss is on that day. The \emph{expected shortfall}\index{expected shortfall}
averages the VaR over the tail beyond $\alpha$,
\begin{equation}\label{eq8:es}
  \mathrm{ES}_\alpha(L) = \frac{1}{1-\alpha}\int_\alpha^1 \mathrm{VaR}_u(L)\dd u ,
\end{equation}
and for a loss with a continuous distribution it equals $\E[L\mid L\ge\mathrm{VaR}_\alpha(L)]$, the
mean loss on the bad days.

\textcite{artzner1999} asked what properties a measure $\rho$ of the risk of a loss should have, and
called a measure \emph{coherent}\index{coherent} when it is monotone (a loss that is larger in every state has larger
risk), translation invariant ($\rho(L+c)=\rho(L)+c$ for a sure amount $c$), positively homogeneous
($\rho(\lambda L)=\lambda\rho(L)$ for $\lambda\ge0$) and subadditive,
\begin{equation}\label{eq8:subadd}
  \rho(L_1+L_2)\le\rho(L_1)+\rho(L_2).
\end{equation}
Subadditivity is the property that matters for a fund of funds or a bank with many desks: the risk of
the whole should not exceed the sum of the risks of its parts, or a risk limit could be met by
splitting a book in two. VaR satisfies the first three but not, in general, the fourth.

\begin{example}{Two loans}\label{ex8:loans}
Two independent loans each lose 100 with probability \saP\% and nothing otherwise. Since
$\saP\%<1\%$, the 99\% VaR of each loan is zero. The probability that at least one defaults is
$1-(1-0.009)^2=\saBoth\%$, which exceeds 1\%, so the 99\% VaR of the pair is 100. Holding both loans
has a VaR of 100 against a sum of zero for the parts, and a lender that splits the pair between two
legal entities appears to remove the risk. The 99\% expected shortfall of one loan averages the worst
1\% of outcomes, which consist of the default with probability \saP\% and a zero loss with
probability 0.1\%, so it is $100\times0.9=\saEsOne$. For the pair the worst 1\% contains the double
default, with probability 0.0081\%, and single defaults, and the ES is \saEsTwo, well below the sum
of 180.
\end{example}

\Cref{fig8:loans} repeats the calculation for every default probability up to 2\%. The VaR of the pair
exceeds the sum of the parts only on a band. The band opens where the chance of at least one default
first exceeds 1\%, at $p=1-\sqrt{0.99}=\saBandLo\%$, and it closes at $p=\saBandHi\%$ inclusive; for any
larger $p$ each loan has a VaR of 100 of its own and the sum is 200. The ES of the pair is at or below the
sum of the two ES figures at every $p$.

\begin{figure}[t]
  \centering
  \includegraphics[width=0.94\textwidth]{ch08/fig10_loans}
  \caption{The 99\% VaR (left) and ES (right) of the two independent loans of \cref{ex8:loans}, each
  losing 100 with probability $p$, for the pair held together and for the sum of the two measured
  apart. In the shaded band, $\saBandLo\%<p\le\saBandHi\%$, the VaR of the pair exceeds the sum; the
  band is shaded in both panels for reference. Below the band the two ES curves coincide. The
  dotted line marks $p=\saP\%$.}
  \label{fig8:loans}
\end{figure}

\begin{proposition}[Expected shortfall is coherent]\label{prop8:es}
Let $L$ have a continuous distribution and a finite mean, and let $q=\mathrm{VaR}_\alpha(L)$. Then
\begin{equation}\label{eq8:esmax}
  \mathrm{ES}_\alpha(L) = \max\bigl\{\E[L\mid A] : \mathbb P(A)=1-\alpha\bigr\},
\end{equation}
the largest mean loss over all events of probability $1-\alpha$, attained by $A=\{L\ge q\}$.
Consequently $\mathrm{ES}_\alpha$ is subadditive, and it is coherent.
\end{proposition}
\begin{proof}
With $U=F_L(L)$, which is uniform on $(0,1)$, the loss is $L=F_L^{-1}(U)$ almost surely, and
$\{L\ge q\}=\{U\ge\alpha\}$ up to a null set. So
$\E[L\,\mathbf 1\{L\ge q\}]=\int_\alpha^1F_L^{-1}(u)\dd u=(1-\alpha)\,\mathrm{ES}_\alpha(L)$, and the
event $\{L\ge q\}$ attains the value in \cref{eq8:esmax}. For any event $A$ with $\mathbb P(A)=1-\alpha$,
\[
  (L-q)\bigl(\mathbf 1\{L\ge q\}-\mathbf 1_A\bigr)\ge0
\]
in every state, since both factors are nonnegative when $L\ge q$ and both are nonpositive when
$L<q$. Taking expectations, and using $\E[\mathbf 1\{L\ge q\}-\mathbf 1_A]=0$, gives
$\E[L\mathbf 1_A]\le\E[L\mathbf 1\{L\ge q\}]$, which is \cref{eq8:esmax}. For two losses with
continuous distributions and continuous sum, let $A^*$ attain the maximum for $L_1+L_2$. Then
$\mathrm{ES}_\alpha(L_1+L_2)=\E[L_1\mid A^*]+\E[L_2\mid A^*]\le\mathrm{ES}_\alpha(L_1)+\mathrm{ES}_\alpha(L_2)$,
because each term is the mean of one loss over an event of probability $1-\alpha$. Monotonicity,
translation invariance and homogeneity follow from \cref{eq8:es}, since $\mathrm{VaR}_u$ has each of
them for every $u$.
\end{proof}

The continuity assumption is only for brevity. \textcite{acerbi2002} show that the average of
\cref{eq8:es} is coherent for every loss with a finite mean, discrete or not, and \cref{ex8:loans}
is a discrete case where it behaves. The market risk standard of
\textcite{bcbs2019} replaced the 99\% VaR of earlier rules with a 97.5\% expected shortfall. The two
levels were chosen to be close for normal losses: $\mathrm{ES}_{0.975}/\mathrm{VaR}_{0.99}$ is
\esRatioNormal{} for any normal law with zero mean. For the daily market returns of this chapter the
ratio of the historical estimates is \esRatioHist, so on this sample the change of measure raises the
risk figure that capital starts from. The capital itself also depends on liquidity horizons, a stressed
period and the rest of the framework.

\FloatBarrier
% =====================================================================================
\section{Four ways to compute a value at risk}
\label{sec8:var}

Over \tdFirst{} to \tdLast{} the US market returned \tdMeanAnn\% a year with a volatility of
\tdVolAnn\%, a daily standard deviation of \tdSd\%. The skewness of daily returns is \tdSkew{} and the
excess kurtosis \tdKurt. The worst day, \tdWorstDate, lost \tdWorst\%, which is \tdWorstSd{} standard
deviations below the mean, and the best, \tdBestDate, gained \tdBest\%. There were \tdBeyondFour{} days
more than four standard deviations from the mean, where a normal law expects \tdBeyondFourNormal, and
\tdBeyondFive{} beyond five, where it expects \tdBeyondFiveNormal. \Cref{fig8:tails} shows the shape.
On a logarithmic scale the normal density falls like a parabola, while the data fall far more slowly
in both tails. A straight line on this scale would mean an exponential tail; whether the tail is a power
law is a question for the log-log axes of \cref{fig8:evt} in \cref{sec8:evt}. A Student $t$ law fitted by
maximum likelihood has only \tdNu{} degrees of freedom.

\begin{figure}[t]
  \centering
  \includegraphics[width=0.94\textwidth]{ch08/fig01_tails}
  \caption{Daily returns of the US market, \tdFirst{} to \tdLast, standardised by their mean and
  standard deviation. Left: histogram on a logarithmic scale with the standard normal density and a
  Student $t$ density fitted by maximum likelihood. Right: sample quantiles against normal quantiles;
  the line is where normal data would lie.}
  \label{fig8:tails}
\end{figure}

There are four common ways to turn such a record into a VaR. The \emph{parametric} method assumes a
law, normal or Student $t$, and uses its quantile with the sample mean and standard deviation.
\emph{Historical simulation} uses the empirical law of past losses: with $n$ observations the VaR at
level $\alpha$ is the $k$th largest loss, $k=\lceil n(1-\alpha)\rceil$, tied losses counted separately, and
the ES is the mean of the $k$ largest. When $n(1-\alpha)$ is not a whole number, the $k$th largest loss is
the quantile of \cref{eq8:var} for the empirical law. When it is a whole number, the quantile of
\cref{eq8:var} is the next smaller loss, the $(k+1)$th largest, and the rule used here deliberately
takes the more conservative $k$th. The mean of the $k$ largest losses is the
ES of the empirical law when $n(1-\alpha)$ is whole, and otherwise a close approximation that gives the
$k$th loss a full weight where the exact value gives it part of one. The \emph{Cornish-Fisher}\index{Cornish-Fisher} method
corrects the normal quantile for the sample skewness $\gamma_3$ and excess kurtosis $\gamma_4$ of returns
\parencite{cornish1938},
\begin{equation}\label{eq8:cf}
  \tilde z = z + \frac{(z^2-1)\gamma_3}{6} + \frac{(z^3-3z)\gamma_4}{24} - \frac{(2z^3-5z)\gamma_3^2}{36},
\end{equation}
with $z=z_{1-\alpha}$, so that $\mathrm{VaR}_\alpha=-\mu-\sigma\tilde z$. Because it needs only four
moments it is popular for hedge funds \parencite{favre2002}. \emph{Monte Carlo} draws from a fitted model, which matters for portfolios of
options and for the GARCH models of \cref{sec8:garch}.

The listing below computes the historical VaR and the exact ES of the empirical law, in which the
boundary loss carries the weight $n(1-\alpha)-(k-1)$. For the \tdN{} days of the sample $n(1-\alpha)$ is
\esTailN{} at 99\% and \esTailS{} at 97.5\%, and the exact ES exceeds the mean of the $k$ largest losses
by \esGapN{} and \esGapS{} percentage points, too little to change any figure in \cref{tab8:var}.
\Needspace*{26\baselineskip}
\lstinputlisting[language=Python]{code/ch08/snippets/hist_var_es.py}

\begin{proposition}[Normal and Student $t$ measures]\label{prop8:normal}
If the return is $\N(\mu,\sigma^2)$ then
\[
  \mathrm{VaR}_\alpha=-\mu+\sigma z_\alpha,\qquad
  \mathrm{ES}_\alpha=-\mu+\sigma\,\frac{\varphi(z_\alpha)}{1-\alpha}.
\]
If the return is $\mu+s\,T$ with $T$ a Student $t$ variable with $\nu>2$ degrees of freedom, density
$g_\nu$ and quantile $t_\alpha$, and $s=\sigma\sqrt{(\nu-2)/\nu}$ so that the standard deviation is
$\sigma$, then $\mathrm{VaR}_\alpha=-\mu+s\,t_\alpha$ and
$\mathrm{ES}_\alpha=-\mu+s\,\dfrac{g_\nu(t_\alpha)}{1-\alpha}\,\dfrac{\nu+t_\alpha^2}{\nu-1}$.
\end{proposition}
\begin{proof}
The loss is $-\mu-\sigma Z$, and $-Z$ has the law of $Z$, so the loss quantile at $\alpha$ is
$-\mu+\sigma z_\alpha$. For the ES, $\E[Z\mathbf 1\{Z>z_\alpha\}]=\int_{z_\alpha}^\infty x\varphi(x)\dd x
=\varphi(z_\alpha)$ because $\varphi'(x)=-x\varphi(x)$; dividing by $1-\alpha$ gives the formula. For
the $t$ law the quantile argument is the same. The density
$g_\nu(x)\propto(1+x^2/\nu)^{-(\nu+1)/2}$ satisfies
$\frac{\dd}{\dd x}\bigl[-\frac{\nu+x^2}{\nu-1}g_\nu(x)\bigr]=x\,g_\nu(x)$, as differentiation shows, so
$\int_{t_\alpha}^\infty x g_\nu(x)\dd x=\frac{\nu+t_\alpha^2}{\nu-1}g_\nu(t_\alpha)$, and the ES follows
in the same way.
\end{proof}

\begin{figure}[t]
  \centering
  \includegraphics[width=0.84\textwidth]{ch08/fig11_cornish}
  \caption{The one day loss quantile of the US market in standard deviations, measured from the mean,
  against the confidence level: the normal law, the historical quantile and the Cornish-Fisher
  expansion \cref{eq8:cf} with the sample skewness \tdSkew{} and excess kurtosis \tdKurt. The horizontal
  axis is stretched so that each step from 90\% to 99\% to 99.9\% takes the same width.}
  \label{fig8:cornish}
\end{figure}

\begin{table}[t]
\centering
\caption{One day VaR and ES of the US market, in percent, from all \tdN{} days, \tdFirst{} to
\tdLast. The $t$ law is fitted by maximum likelihood with free location and scale; the
Cornish-Fisher method uses the sample skewness and kurtosis and has no ES of its own.}
\label{tab8:var}
\small
\begin{tabular}{lrrrr}
\toprule
 & \multicolumn{2}{c}{99\%} & \multicolumn{2}{c}{97.5\%}\\
\cmidrule(lr){2-3}\cmidrule(lr){4-5}
 & VaR & ES & VaR & ES\\
\midrule
Normal & \vrNorN & \esNorN & \vrNorS & \esNorS\\
Student $t$ & \vrTtN & \esTtN & \vrTtS & \esTtS\\
Cornish-Fisher & \vrCfN & & \vrCfS & \\
Historical & \vrHsN & \esHsN & \vrHsS & \esHsS\\
\bottomrule
\end{tabular}
\end{table}

\Cref{tab8:var} compares the four methods on the whole sample. The normal VaR at 99\% is \vrNorN\%
against a historical \vrHsN\%, and the normal ES is \esNorN\% against \esHsN\%: the normal law
understates the 99\% loss by a sixth and the mean loss beyond it by more than a third. The $t$ law
comes close to the historical VaR at both levels and overstates the ES somewhat, because three degrees of
freedom give a tail heavier than the data's. The Cornish-Fisher VaR of \vrCfN\% is more than twice the
historical figure. \Cref{eq8:cf} is a series expansion that is accurate when skewness and kurtosis are
small, and with an excess kurtosis of \tdKurt{} the cubic term dominates; at 99.9\% the corrected
loss quantile is \cfLossNine{} standard deviations from the mean, against \cfHsThree{} in the data, an
absurd value for a loss expected about once in four years. The expansion also stops being monotone
in $\alpha$ when the moments are large, so it can rank two levels in the wrong order.
\Cref{fig8:cornish} shows both failures for these moments. The corrected VaR falls as the level rises
from 50\% to \cfDownHi\%, it is negative, a gain, up to \cfNegHi\%, and beyond 99\% it climbs far above
the historical quantile. It is a tool for
the mild non-normality of monthly fund returns, and even there the moments it needs are poorly
estimated from a short record.

A daily VaR is often turned into a ten day one by the square root of time\index{square root of time}. If daily returns are
independent and identically distributed, the variance of a ten day sum is ten times the daily
variance, so for normal returns $\mathrm{VaR}$ over ten days is $\sqrt{10}$ times the daily VaR when the
mean is ignored. For heavy-tailed returns the rule is only approximate even under independence: with a
finite variance, the law of a ten day sum is closer to normal than the law of one day, so its quantile
is not $\sqrt{10}$ times the daily one. On this sample $\sqrt{10}$ times the daily
historical VaR of \srOne\% is \srScaled\%, while the 99\% quantile of overlapping ten day losses is
\srTen\%; for the ES the figures are \srScaledEs\% and \srTenEs\%. The rule errs on the safe side
here, but it rests on independence, and the next section shows that daily returns are uncorrelated
without being independent.

\FloatBarrier
% =====================================================================================
\section{Volatility that comes in clusters}
\label{sec8:garch}

The lag one autocorrelation of daily market returns is \acRet, too small to trade on, but that of
squared returns is \acSq{} and that of absolute returns \acAbs. Large moves follow large moves of
either sign, and the effect lasts: in the left panel of \cref{fig8:clusters} every one of the first
\acLags{} autocorrelations of squared and of absolute returns lies above the band $\pm2/\sqrt n$ that
independent data would stay inside about 95\% of the time, and at lag \acLags{} the autocorrelation of
absolute returns is still \acAbsHundred. Those of returns stay close to the band, although \acRetOut{} of
the \acLags{} fall outside it where independence would put about five; when volatility clusters, the band
computed for independent data is itself too narrow. A measure computed from the whole sample, or from a long window, therefore averages calm
and turbulent years and is wrong in both. \textcite{engle1982} modelled the conditional variance as a
function of past squared shocks, and the generalised form of \textcite{bollerslev1986} is the one in
general use:
\begin{equation}\label{eq8:garch}
  r_t=\mu+\varepsilon_t,\qquad \varepsilon_t=\sigma_t z_t,\qquad
  \sigma_t^2=\omega+a\,\varepsilon_{t-1}^2+b\,\sigma_{t-1}^2,
\end{equation}
with $z_t$ independent, mean zero and variance one, and $\omega>0$, $a\ge0$, $b\ge0$. The letters $a$
and $b$ are used here so that $\alpha$ can remain a confidence level. When $a+b<1$, writing
$\omega=(1-a-b)\bar\sigma^2$ shows that tomorrow's variance is a weighted average of the long run variance
$\bar\sigma^2$, today's squared shock and today's variance, with weights $1-a-b$, $a$ and $b$.

\begin{proposition}[Moments and forecasts of GARCH(1,1)]\label{prop8:garch}
Let $p=a+b<1$ and $\bar\sigma^2=\omega/(1-p)$. In the stationary solution of \cref{eq8:garch}:
\begin{enumerate}[label=(\roman*)]
  \item the unconditional variance of $\varepsilon_t$ is $\bar\sigma^2$;
  \item the forecast made at $t$ of the variance $k\ge1$ days ahead is
  $\E_t[\sigma_{t+k}^2]=\bar\sigma^2+p^{\,k-1}(\sigma_{t+1}^2-\bar\sigma^2)$, so a shock to the variance
  halves in $\ln\tfrac12/\ln p$ days;
  \item if the $z_t$ are normal and $p^2+2a^2<1$, the kurtosis of $\varepsilon_t$ is
  $3(1-p^2)/(1-p^2-2a^2)$, which exceeds 3 whenever $a>0$.
\end{enumerate}
\end{proposition}
\begin{proof}
Since $z_t$ is independent of the past with variance one, $\E_{t-1}[\varepsilon_t^2]=\sigma_t^2$. For
$k\ge2$, $\E_t[\sigma_{t+k}^2]=\omega+a\E_t[\varepsilon_{t+k-1}^2]+b\E_t[\sigma_{t+k-1}^2]
=\omega+p\,\E_t[\sigma_{t+k-1}^2]$. Subtracting $\bar\sigma^2=\omega+p\bar\sigma^2$ gives
$\E_t[\sigma_{t+k}^2]-\bar\sigma^2=p\,(\E_t[\sigma_{t+k-1}^2]-\bar\sigma^2)$, and iterating from $k=1$
proves (ii). Taking unconditional expectations in the recursion gives $m=\omega+pm$ for
$m=\E[\sigma_t^2]=\E[\varepsilon_t^2]$, which is (i). For (iii), $\E[\varepsilon_t^4]=3\E[\sigma_t^4]$ for
normal $z_t$, and squaring the recursion,
\[
  \E[\sigma_t^4]=\omega^2+2\omega p\,m+(3a^2+2ab+b^2)\,\E[\sigma_{t-1}^4],
\]
because $\E[z^4]=3$. Here $3a^2+2ab+b^2=p^2+2a^2$, so in the stationary solution
$\E[\sigma^4](1-p^2-2a^2)=\omega^2(1+p)/(1-p)$, and dividing $3\E[\sigma^4]$ by $m^2=\omega^2/(1-p)^2$
gives the kurtosis.
\end{proof}

When $p^2+2a^2\ge1$ the recursion for the fourth moment has no finite solution, and the model, though
stationary, has infinite kurtosis. The RiskMetrics rule \parencite{riskmetrics1996} sets
$\sigma_t^2=\lambda\sigma_{t-1}^2+(1-\lambda)r_{t-1}^2$ with $\lambda=0.94$ for daily data. It is the
case $\omega=0$, $a=1-\lambda$, $b=\lambda$ of \cref{eq8:garch} with zero mean, so $p=1$, which (ii)
excludes, and there is no long run variance. Its forecasts follow directly: $\E_t[r_{t+j}^2]=\E_t[\sigma_{t+j}^2]$
for $j\ge1$, so $\E_t[\sigma_{t+k}^2]=\lambda\E_t[\sigma_{t+k-1}^2]+(1-\lambda)\E_t[\sigma_{t+k-1}^2]=\sigma_{t+1}^2$, and the
forecast of the variance stays at tomorrow's level at every horizon.

Fitted by maximum likelihood to all \tdN{} days, with the variance started at the sample variance,
the normal GARCH model has $a=\gaAlphaN$ and $b=\gaBetaN$, a persistence of \gaPersN, a half life of
\gaHalfN{} trading days for a shock to the variance and a long run volatility of \gaLrVolN\% a year.
With Student $t$ innovations the estimates are $a=\gaAlphaT$ and $b=\gaBetaT$ with $\nu=\gaNu$, and
the likelihood ratio statistic against the normal model is \gaLR{} for one extra parameter. The $t$
model is more persistent, at \gaPersT, which puts its long run volatility at \gaLrVolT\%. That figure
is far from the sample's \tdVolAnn\% and shows how poorly $\bar\sigma^2=\omega/(1-p)$ is determined when
$p$ is close to one. For the normal fit $p^2+2a^2=\gaFourth$, just above one, so the fitted model has
no finite fourth moment, which is consistent with a sample kurtosis that is large and unstable.
The R listing below fits the normal model to the same returns in percent. It maps the parameters so
that $\omega>0$, $a\ge0$, $b\ge0$ and $a+b<1$ hold for any input, refines a simplex search with a
quasi-Newton step, and prints the estimates above, which a test compares with the Python fit.
\Needspace*{19\baselineskip}
\lstinputlisting[language=R]{code/ch08/snippets/garch_mle.R}

The model removes most of the lag one clustering in squared returns. The standardised residuals $z_t$
have a lag one autocorrelation of squares of \gaResAcSq{} against \acSq{} for the raw returns, and excess
kurtosis of \gaResKurt{} against \tdKurt. What remains is still not normal: the skewness of the residuals
is \gaResSkew, the same as that of the raw returns to two decimals, and their 1\% quantile is $-\gaResQ$, against $-2.33$ for a standard normal. That gap is
the reason for the $t$ innovations and for the filtered historical simulation of
\cref{sec8:case}. The right panel of \cref{fig8:clusters} sets the residuals against the $t$ law with the
\gaNu{} degrees of freedom of the second fit. The $z_t$ of \cref{eq8:garch} have variance one, so the
reference must too: the $t$ quantiles are multiplied by $\sqrt{(\nu-2)/\nu}$, which moves the 1\%
quantile from \gaTQRaw{} to \gaTQ. Even so, the 0.1\% quantile of the residuals is \gaResQThou,
against \gaTQThou{} for the rescaled $t$ (computed with the unrounded $\nu=\gaNuFine$) and \gaNQThou{} for the normal, while the right tail of the
residuals is lighter than that of the $t$. The EWMA variance is close to the GARCH variance, with a correlation of \ewCorr{}
between the two volatility series, but its normal log likelihood, computed on the same days and returns
in the same units, is lower by \ewLLdiff.

\begin{figure}[t]
  \centering
  \includegraphics[width=0.94\textwidth]{ch08/fig12_clusters}
  \caption{Left: autocorrelations of daily US market returns, of their squares and of their absolute
  values at lags 1 to \acLags, with the band $\pm2/\sqrt n$ (shaded). Right: quantiles of the
  standardised residuals of the normal GARCH(1,1) fit against those of the standard normal and of the
  Student $t$ law with \gaNu{} degrees of freedom rescaled to variance one. The line is where the
  residuals would lie if the reference law were right.}
  \label{fig8:clusters}
\end{figure}

\begin{figure}[t]
  \centering
  \includegraphics[width=0.94\textwidth]{ch08/fig02_garch}
  \caption{Volatility of the US market from the normal GARCH(1,1) model fitted to the days from \tdFirst{}
  to \tdLast, in percent a year, with $\sqrt{252}\,|r_t-\mu|$, the absolute deviation of each daily
  return from the mean in annual units, behind it.}
  \label{fig8:garch}
\end{figure}

\Cref{fig8:garch} shows the fitted volatility. Its median is \volMedian\% a year, its lowest value
\volMin\% and its highest \volMax\%, on \volMaxDate. \Cref{fig8:forecast} turns
\cref{prop8:garch}(ii) into forecasts of the average volatility over the next $k$ trading days from
four dates. Here the average volatility over $k$ days is the square root of the average variance,
$\bigl(\tfrac{252}{k}\sum_{j=1}^k\E[\sigma_{t+j}^2]\bigr)^{1/2}$, the annualised daily-equivalent standard deviation of the
$k$-day cumulative return, which is the root of the summed variances because the model's returns are
uncorrelated, and not the average of the daily volatilities. Each forecast uses data to the close of
the trading day before its date, and the realised figure is $\sqrt{252}$ times the standard deviation of
the 252 daily returns from that date on. For \fcKaDate, the day after the crash, the model's volatility
was \fcKaNow\% and its forecast for the next year \fcKaYear\%; the volatility realised over that year
was \fcKaReal\%. For \fcKbDate{} the figures are \fcKbNow, \fcKbYear{} and \fcKbReal, and for
\fcKcDate, whose forecast uses data to the close of 13 March and so precedes that day's fall, they are
\fcKcNow, \fcKcYear{} and \fcKcReal. In the calm of \fcKdDate{} the model forecast \fcKdYear\% from a
starting level of \fcKdNow\% and \fcKdReal\% was realised. Volatility after each crisis fell faster
than a persistence of \gaPersN{} implies, so after a spike the one year forecasts were well above the
volatility that followed; four dates say nothing about how well the one day forecasts are calibrated,
which is the subject of \cref{sec8:case}. A risk manager who scales a one day GARCH VaR to a year
by \cref{prop8:garch}(ii) carries the persistence estimate, and its error, into the answer.

\begin{figure}[t]
  \centering
  \includegraphics[width=0.84\textwidth]{ch08/fig03_forecast}
  \caption{Forecasts from the normal GARCH(1,1) model of the average volatility over the next $k$
  trading days, the square root of 252 times the mean forecast daily variance over those days. Each is
  made with data to the close of the trading day before the date shown, which is the first day forecast.
  The dotted line is the long run volatility of the fitted model.}
  \label{fig8:forecast}
\end{figure}

\FloatBarrier
% =====================================================================================
\section{The far tail}
\label{sec8:evt}

Historical simulation has nothing to say beyond the largest loss in the sample, and parametric laws
fitted to the whole distribution are fitted mostly to its centre. Extreme value theory models the
tail alone. For a loss $L$ and a high threshold $u$, the excess $L-u$ given $L>u$ is, for a wide class
of laws, close to a generalised Pareto law with distribution function
\begin{equation}\label{eq8:gpd}
  G_{\xi,\beta}(y)=1-\Bigl(1+\frac{\xi y}{\beta}\Bigr)^{-1/\xi},\qquad y\ge0,\quad 1+\xi y/\beta>0,
\end{equation}
with scale $\beta>0$, so that for $\xi<0$ the law ends at $-\beta/\xi$, and with $G_{0,\beta}(y)=1-e^{-y/\beta}$,
when $u$ is large \parencite{balkema1974,pickands1975}. The shape $\xi$ sets the weight of the tail:
for $\xi>0$ the tail falls like a power with index $1/\xi$, so moments of order $1/\xi$ and higher are
infinite, and for the normal law the limit has $\xi=0$. \textcite{embrechts1997} give the theory and
\textcite{mcneil2000} apply it to the residuals of a GARCH model.

\begin{proposition}[Value at risk and expected shortfall from a fitted tail]\label{prop8:gpd}
Suppose that $n$ losses include $n_u$ above $u$ and that for $y\ge0$,
$\mathbb P(L>u+y)=\frac{n_u}{n}\bigl(1-G_{\xi,\beta}(y)\bigr)$ with $0<\xi<1$. Then for $\alpha>1-n_u/n$,
\[
  \mathrm{VaR}_\alpha=u+\frac{\beta}{\xi}\Bigl[\Bigl(\frac{n}{n_u}(1-\alpha)\Bigr)^{-\xi}-1\Bigr],\qquad
  \mathrm{ES}_\alpha=\frac{\mathrm{VaR}_\alpha+\beta-\xi u}{1-\xi}.
\]
\end{proposition}
\begin{proof}
Setting $\frac{n_u}{n}(1+\xi y/\beta)^{-1/\xi}=1-\alpha$ and solving for $y$ gives the VaR as $u+y$.
For $v\ge u$ the excess of a generalised Pareto variable $Y$ over $w=v-u$ satisfies
\[
  \mathbb P(Y>w+x\mid Y>w)=\Bigl(\frac{\beta+\xi(w+x)}{\beta+\xi w}\Bigr)^{-1/\xi}
  =\Bigl(1+\frac{\xi x}{\beta+\xi w}\Bigr)^{-1/\xi},
\]
so it is generalised Pareto with the same $\xi$ and scale $\beta+\xi w$, and its mean is
$(\beta+\xi w)/(1-\xi)$ for $\xi<1$. With $v=\mathrm{VaR}_\alpha$ the ES is
$v+(\beta+\xi(v-u))/(1-\xi)=(v+\beta-\xi u)/(1-\xi)$.
\end{proof}

With the threshold at the 95th percentile of daily losses, $u=\evU\%$, there are \evNu{} exceedances,
and the maximum likelihood fit gives $\xi=\evXi$ with a standard error of about \evXiSe, and
$\beta=\evBeta\%$. The tail index $1/\xi$ is \evAlphaTail, just below four, so the fitted Pareto tail, taken
literally, has no finite fourth moment, which is what the GARCH fit suggested too. The estimate does
not settle the question: two standard errors either side of $\xi$ put the tail index between
\evAlphaLo{} and \evAlphaHi, which includes laws with a finite fourth moment and laws without one. \Cref{fig8:evt} compares the fitted
tail with the data. At 99\% the fitted VaR is \evVarN\% against a historical \evHsVarN\%, and at
99.9\% it is \evVarT\% against \evHsVarT\%, with an ES of \evEsT\%. At 99.99\%, a loss expected once in
forty years of trading days, the fitted VaR is \evVarF\% and the ES \evEsF\%; the historical estimate,
\evHsVarF\%, is the second largest loss in the sample, and the normal law gives \evNorVarF\%. Under the
fitted tail, taken to be stationary and extrapolated beyond the data, a loss as large as that of 19
October 1987 has a return period of about \evYears{} years. Under the normal law it is
\evNormalYears{} years. Both are statements about models, not frequencies that the data could check.

\begin{figure}[t]
  \centering
  \includegraphics[width=0.94\textwidth]{ch08/fig04_evt}
  \caption{Left: the probability that a daily loss of the US market exceeds a given size, from the
  data above the 95th percentile, the fitted generalised Pareto tail and the normal law, both axes
  logarithmic. Right: the Hill estimate of the tail index $1/\xi$ against the number $k$ of largest
  losses used; the dashed line is the index of the generalised Pareto fit.}
  \label{fig8:evt}
\end{figure}

\begin{figure}[!tb]
  \centering
  \includegraphics[width=0.8\textwidth]{ch08/fig13_threshold}
  \caption{Left: the shape $\xi$ of the generalised Pareto law fitted to daily US market losses above
  each threshold from the 90th to the 99th percentile, with plus and minus two standard errors,
  $(1+\xi)/\sqrt{n_u}$, at each threshold taken alone (shaded); the dotted line is the threshold used in
  the text. Right: the mean excess over a threshold, which is linear in the threshold for a generalised
  Pareto tail, with the line implied by the fit at the 95th percentile.}
  \label{fig8:threshold}
\end{figure}

Two cautions. The first is that the estimate of $\xi$ depends on the threshold. The Hill estimator\index{Hill estimator} \parencite{hill1975},
which assumes an exact power tail above the $k$th largest loss, puts the tail index $1/\xi$ at
\evHillTwoHundred{} with $k=200$ and \evHillEight{} with $k=800$, below the generalised Pareto value of
\evAlphaTail, and the right panel of \cref{fig8:evt} drifts with $k$ instead
of settling. The generalised Pareto fit drifts too (\cref{fig8:threshold}): as the threshold moves from
the 90th to the 99th percentile, the fitted $\xi$ rises from \evXiMin{} to a peak of \evXiMax{} at the
\evXiMaxQ th percentile, and at the 99th percentile, with only \evNuTop{} exceedances, its standard error
is \evXiSeTop. The
shaded region in that figure is plus and minus two standard errors at each threshold taken alone. It is
not a confidence band for the curve as a whole: the fits share most of their data, and the chance that
the whole curve lies inside a set of pointwise intervals is lower than the chance for any one point. The second caution is that even at a
fixed threshold the fitted shape can be biased when the threshold is not yet far enough into the tail,
and a simulation shows it. For \evSimM{} samples of the same length from a
Student $t$ law with four degrees of freedom, whose true $\xi$ is $0.25$, the same procedure gives
$\xi$ with a mean of \evSimXi{} and a standard deviation of \evSimXiSd: the 95th percentile is not yet
far enough into the tail of this law. Yet its 99.9\% VaR, measured relative to the true value, has a
mean error of only \evSimBias\% and a standard deviation across samples of \evSimSd\%, against a
standard deviation of \evSimHsSd\% for the historical quantile. In this design the fitted tail
estimates the high quantile with a smaller spread than the sample does, although its shape parameter
is biased.


\FloatBarrier
% =====================================================================================
\section{Testing a risk model against what happened}
\label{sec8:backtest}

A VaR model makes a forecast every day, and every day the forecast can be checked. Define the
\emph{hit} on day $t$ as $I_t=\mathbf 1\{L_t>\mathrm{VaR}_t\}$, where $\mathrm{VaR}_t$ was computed
with information up to the close of day $t-1$. If the model is right, the hits have a simple law.

\begin{proposition}[Hits of a correct model]\label{prop8:hits}
Let $\mathcal F_{t-1}$ be the information at the close of day $t-1$, and suppose that given
$\mathcal F_{t-1}$ the loss $L_t$ has a continuous distribution whose $\alpha$ quantile is
$\mathrm{VaR}_t$. Then $I_1,\dots,I_T$ are independent Bernoulli variables with success probability
$1-\alpha$, and the number of exceptions $x=\sum_t I_t$ is binomial with parameters $T$ and $1-\alpha$.
\end{proposition}
\begin{proof}
By assumption $\mathbb P(I_t=1\mid\mathcal F_{t-1})=1-\alpha$, a constant. The hits $I_1,\dots,I_{t-1}$ are
known at $t-1$, so for any values $i_1,\dots,i_t$,
\[
  \mathbb P(I_1=i_1,\dots,I_t=i_t)=\E\bigl[\mathbf 1\{I_1=i_1,\dots,I_{t-1}=i_{t-1}\}\,
  \mathbb P(I_t=i_t\mid\mathcal F_{t-1})\bigr],
\]
and the conditional probability is $(1-\alpha)^{i_t}\alpha^{1-i_t}$ whatever happened before. Induction
on $t$ gives the product of $T$ Bernoulli probabilities.
\end{proof}

The proposition gives two things to test. \textcite{kupiec1995} tests the count. With
$\pi_0=1-\alpha$ and $\hat\pi=x/T$, the statistic
\begin{equation}\label{eq8:kupiec}
  \mathrm{LR}_{\mathrm{uc}}=-2\ln\frac{(1-\pi_0)^{T-x}\pi_0^{\,x}}{(1-\hat\pi)^{T-x}\hat\pi^{\,x}}
\end{equation}
is approximately $\chi^2_1$ under the null. \textcite{christoffersen1998} tests independence against a
first order Markov chain in which the probability of a hit depends on whether there was a hit the day
before. With $n_{ij}$ the number of days in state $j$ that follow a day in state $i$,
$\hat\pi_{01}=n_{01}/(n_{00}+n_{01})$ and $\hat\pi_{11}=n_{11}/(n_{10}+n_{11})$, the statistic
compares the likelihood of the chain with that of independent hits at the overall rate and is also
approximately $\chi^2_1$. A model that is right on average but whose exceptions bunch in crises can
pass the first test and fail the second.

The listing below computes both statistics from a series of hits. It uses the convention $0\ln0=0$, so
a record in which no exception follows another, $n_{11}=0$, gives $\hat\pi_{11}=0$ and a finite
statistic, and so does a record with no exceptions at all. Applied to the hit series of
\cref{sec8:case} it reproduces the $p$ values for the US market in \cref{tab8:case}, which a test checks.
\Needspace*{22\baselineskip}
\lstinputlisting[language=Python,linerange={1-19}]{code/ch08/snippets/lr_tests.py}
\Needspace*{18\baselineskip}
\lstinputlisting[language=Python,linerange={22-37}]{code/ch08/snippets/lr_tests.py}
\Needspace*{9\baselineskip}
\lstinputlisting[language=Python,linerange={40-45}]{code/ch08/snippets/lr_tests.py}

Regulators count. The framework of \textcite{bcbs1996}, the original rules for internal models,
sorts a bank's last 250 days of exceptions at 99\% into three zones: up to four is green, five to nine
is yellow and ten or more is red. In those rules the multiplier on the bank's capital rises from three
toward four through the yellow zone and is four in the red. The current standard of \textcite{bcbs2019}
keeps a count of exceptions over 250 days with three zones but sets capital differently, so the
multipliers here describe the 1996 framework only. For a correct model the probabilities follow from the binomial law: \znGreen\% green,
\znYellow\% yellow and \znRed\% red. The design accepts a high chance of a false yellow to keep a false
red almost impossible, and the price is weak power. At a true exception rate of 2\%, twice the
nominal, the model is still green with probability \znGreenTwo\% and red with probability
\znRedTwo\%; at 3\% the figures are \znGreenThree\% and \znRedThree\%.

The Kupiec\index{Kupiec} test at the 5\% level does no better over one year. It accepts any count from \kpLo{} to
\kpHi{} out of 250; six exceptions give $\mathrm{LR}_{\mathrm{uc}}=\kpSixLR$ and a $p$ value of \kpSixP.
Because the count is discrete, rejecting whenever the chi-square $p$ value is below 5\% rejects a
correct model over 250 days with probability \pwSizeOne\%, not 5\%. \Cref{fig8:power}
shows the exact rejection probability as the sample grows (a year here is 250 trading days, as in
the Basel\index{Basel} count; the curves are jagged because the acceptance region moves by whole exceptions, so
power can dip when a year is added). A model with a true rate of 2\% is rejected
with probability \pwOneTwo\% after one year, and a model at 1.5\% with probability \pwOneFifteen\%,
\pwFiveFifteen\% after five years and \pwTenFifteen\% after ten. Power of 80\% needs \pwYearsTwo{}
years of data at 2\% and \pwYearsFifteen{} years at 1.5\%. A year of exceptions can expose a badly wrong
model and very little else.

\begin{figure}[t]
  \centering
  \includegraphics[width=0.94\textwidth]{ch08/fig05_power}
  \caption{Left: exact probability that the Kupiec test at the 5\% level rejects a 99\% VaR model
  whose true exception rate is 2\% or 1.5\%, against years of daily data; the dotted line marks a
  power of 80\%. Right: probabilities of the three Basel zones over 250 days against the true
  exception rate.}
  \label{fig8:power}
\end{figure}

Expected shortfall is harder to test because it is an average over the tail, so a single day says
little about it. It is also not \emph{elicitable}\index{elicitable} \parencite{gneiting2011}: no scoring function of an ES forecast and the outcome has the
true ES as the minimiser of its expected value, a point that \textcite{acerbi2014} discuss before proposing tests that
use the size of the losses on exception days. The pair of VaR and ES is jointly elicitable
\parencite{fissler2016}, which is what makes comparisons of competing forecasts of the pair possible. One of the Acerbi-Sz\'ekely statistics, in the sign convention used here, is
\begin{equation}\label{eq8:z2}
  Z_2=\frac{1}{T(1-\alpha)}\sum_{t=1}^{T}\frac{L_t I_t}{\mathrm{ES}_t}-1,
\end{equation}
with $I_t$ the hit against the VaR at the same level $\alpha$. It is a statistic, not a complete test:
its law under the null must be worked out for the loss distribution in question.

\begin{proposition}[The shortfall test is centred]\label{prop8:z2}
If given $\mathcal F_{t-1}$ the loss $L_t$ has a continuous distribution and $\mathrm{VaR}_t$ and
$\mathrm{ES}_t$ are its true $\alpha$ quantile and expected shortfall, then $\E[Z_2]=0$. If the model's
VaR is right but its ES is too low on every day, $\E[Z_2]>0$.
\end{proposition}
\begin{proof}
For a continuous law $\mathbb P(I_t=1\mid\mathcal F_{t-1})=1-\alpha$ and
$\mathrm{ES}_t=\E[L_t\mid L_t>\mathrm{VaR}_t,\mathcal F_{t-1}]$, so
$\E[L_tI_t\mid\mathcal F_{t-1}]=(1-\alpha)\mathrm{ES}_t$. Dividing by the known $\mathrm{ES}_t$ and taking
expectations gives $\E[L_tI_t/\mathrm{ES}_t]=1-\alpha$ for every $t$, and $\E[Z_2]=0$. If the model
reports $\widetilde{\mathrm{ES}}_t<\mathrm{ES}_t$ with the true VaR, each term has expectation
$(1-\alpha)\mathrm{ES}_t/\widetilde{\mathrm{ES}}_t>1-\alpha$.
\end{proof}

The law of $Z_2$ under the null depends on the loss distribution. For independent losses with fixed,
correctly specified forecasts of VaR and ES and a finite variance of $X_t$, the ordinary central limit
theorem gives an approximation. Write $X_t=L_tI_t/\mathrm{ES}_t$. Then $\E X_t=1-\alpha$ for a continuous
law, $Z_2=\bar X/(1-\alpha)-1$ has mean zero, and by the central limit theorem its standard deviation is
$\sqrt{\operatorname{Var}X_t/T}/(1-\alpha)$, where $\operatorname{Var}X_t$ follows from the law by one
integral. For \ztT{} days of Student $t$ losses with five degrees of freedom at 97.5\%, the standard
deviation is \ztSdAn{} and the one-sided 5\% critical value \ztCritAn. A simulation of \ztM{} such
samples, in \ztStreams{} independent streams, gives a mean of \ztMean, a standard deviation of \ztSd{}
and a 95th percentile of \ztCrit{} with a standard error of about \ztCritSeAn, from the asymptotic
variance of a sample quantile with the normal density at the percentile (the spread of the
percentile across the streams gives \ztCritSe), so the normal approximation is
adequate at this length. For a zero-location scale family the ratio $L_t/\mathrm{ES}_t$ does not depend on
the scale. With a nonzero location it also depends on the location relative to the scale, so a
critical value calculated for one model need not apply to another; the daily means here are small
against the volatility, and the same calculation, run with the location set to zero, gives a
reference critical value for each model from the law it assumes, which \cref{sec8:case} uses. They are
reference values: forecasts that move with a conditional volatility and parameters estimated from past
returns add dependence and estimation error, and a null law that allows for both needs a martingale
limit theorem or a simulation of the whole forecasting procedure. Too many exceptions also push
$Z_2$ up, so the test responds to a wrong VaR as well as a wrong ES; that is a feature, since the
shortfall depends on both.

The Octave listing below does the integral for the $t$ law and prints \ztSdAn{} and \ztCritAn. The
variable $X_t$ is zero on a fraction $\alpha$ of days, and those zeros belong in its variance:
$\operatorname{Var}X_t=\E X_t^2-(1-\alpha)^2$, where $\E X_t^2$ integrates the squared scaled loss
over the tail alone because the zeros add nothing to it. The variance of the loss beyond the VaR,
taken by itself, leaves out the days without an exception and gives a different figure.
\Needspace*{17\baselineskip}
\lstinputlisting[language=Octave]{code/ch08/snippets/es_null.m}

\FloatBarrier
% =====================================================================================
\section{Return per unit of risk}
\label{sec8:ratios}

A risk measure answers how much could be lost. An investor choosing between managers also wants to
know how much return each unit of risk bought. Write $R$ for a portfolio's return in excess of the
risk free rate, with mean $\mu$, volatility $\sigma$ and beta $\beta$ on the market excess return
$R_m$, and $e$ for the residual of the regression of $R$ on $R_m$, with intercept $a_J$. The common
measures are:
\begin{itemize}
  \item the \emph{Sharpe ratio}\index{Sharpe ratio} $\mu/\sigma$ \parencite{sharpe1966,sharpe1994};
  \item the \emph{Treynor ratio}\index{Treynor ratio} $\mu/\beta$, return per unit of market risk \parencite{treynor1965};
  \item \emph{Jensen's alpha}\index{Jensen's alpha} $a_J=\mu-\beta\mu_m$ \parencite{jensen1968};
  \item the \emph{appraisal} or \emph{information ratio}\index{information ratio} $a_J/\sigma_e$, alpha per unit of the
        residual risk $\sigma_e$, which is the tracking error against the beta adjusted market;
  \item $M^2=(\mu/\sigma)\,\sigma_m$, the excess return of the portfolio levered or delevered to the
        market's volatility \parencite{modigliani1997};
  \item the \emph{Sortino ratio}\index{Sortino ratio} $\mu/\sqrt{\E[\min(R,0)^2]}$, which counts only returns below a
        target, here zero \parencite{sortino1994};
  \item the \emph{Omega ratio}\index{Omega ratio} $\Omega(\tau)=\E[(R-\tau)^+]/\E[(\tau-R)^+]$ at a threshold $\tau$
        \parencite{keating2002}.
\end{itemize}
In practice the information ratio is also computed from returns in excess of a benchmark, with no
beta adjustment; the two coincide when beta is one.

\begin{proposition}[What $M^2$ and Omega measure]\label{prop8:m2}
\begin{enumerate}[label=(\roman*)]
  \item A position of $w=\sigma_m/\sigma$ in the portfolio and $1-w$ in the risk free asset has the
  market's volatility and excess return $M^2$, so ranking by $M^2$ is ranking by the Sharpe ratio.
  \item For any threshold $\tau$, $\Omega(\tau)=1+(\mu-\tau)/\E[(\tau-R)^+]$. So $\Omega(\tau)>1$ exactly
  when $\mu>\tau$, and Omega depends on the law of $R$ only through its mean and its first lower partial
  moment $\E[(\tau-R)^+]$.
\end{enumerate}
\end{proposition}
\begin{proof}
(i) The position has excess return $wR$, with mean $w\mu=\mu\sigma_m/\sigma$ and volatility
$w\sigma=\sigma_m$. Since $\sigma_m$ is common to every portfolio, $M^2$ orders portfolios as $\mu/\sigma$
does. (ii) $(R-\tau)^+-(\tau-R)^+=R-\tau$, so $\E[(R-\tau)^+]=\mu-\tau+\E[(\tau-R)^+]$, and dividing by
$\E[(\tau-R)^+]>0$ gives the formula.
\end{proof}

For the US market from \prFirst{} to \prLast, \prT{} months, the Sharpe ratio is \prSR{} a year, or
\prSRm{} a month. That figure is an estimate, and the next result gives its standard error.

\begin{proposition}[Standard error of the Sharpe ratio]\label{prop8:sharpese}
Let returns be independent and identically distributed with Sharpe ratio $\theta$, skewness $\gamma_3$,
excess kurtosis $\gamma_4$ and a finite fourth moment, and let $\hat\theta$ be the sample mean over the
sample standard deviation of $T$ returns. Then $\sqrt T(\hat\theta-\theta)$ is asymptotically normal
with variance
\begin{equation}\label{eq8:sharpese}
  1+\tfrac12\theta^2-\gamma_3\theta+\tfrac14\gamma_4\theta^2 .
\end{equation}
\end{proposition}
\begin{proof}
By the central limit theorem the sample mean $\hat\mu$ and variance $\hat v$ are jointly asymptotically
normal with $T\Var(\hat\mu)\to\sigma^2$, $T\Var(\hat v)\to\mu_4-\sigma^4=\sigma^4(2+\gamma_4)$ and
$T\Cov(\hat\mu,\hat v)\to\mu_3=\gamma_3\sigma^3$, where $\mu_3$ and $\mu_4$ are central moments. The map
$g(\mu,v)=\mu/\sqrt v$ has partial derivatives $1/\sigma$ and $-\mu/(2\sigma^3)$, so by the delta
method the asymptotic variance of $\sqrt T\hat\theta$ is
\[
  \frac{\sigma^2}{\sigma^2}+\frac{\mu^2}{4\sigma^6}\,\sigma^4(2+\gamma_4)
  -2\,\frac{1}{\sigma}\,\frac{\mu}{2\sigma^3}\,\gamma_3\sigma^3
  =1+\frac{\theta^2(2+\gamma_4)}{4}-\gamma_3\theta,
\]
which is \cref{eq8:sharpese}.
\end{proof}

For independent normal returns the skewness and excess kurtosis vanish, leaving $1+\theta^2/2$, the
case treated by \textcite{lo2002}, whose adjustment for serially dependent returns appears below;
\textcite{opdyke2007} gives the general form and extends it to dependent returns. For a positive
Sharpe ratio negative skewness widens the interval, and fat tails widen it whatever the sign of the
ratio. For the market, with monthly skewness \prSkew{} and excess kurtosis \prKurt, the annualised
standard error is \prSeIid{} under the normal law and \prSeNn{} in general. Over \prT{} months that is
precise enough; over ten years, at the same ratio, the standard error is \prSeTen, so a two standard
error interval for a manager with the market's Sharpe ratio would include zero. Comparing two managers
would also need the correlation of their returns. \Cref{fig8:sharpese} turns the question round: how
many years of independent monthly returns bring the standard error down to half the true ratio, so
that an estimate equal to the true ratio lies two standard errors above zero? At the market's ratio the
answer is about \prYearsTwo{} years. For a ratio of one it is \prYearsOne{} years, or \prYearsOneAlt{}
if the returns have skewness \prAltSkew{} and excess kurtosis \prAltKurt.

\begin{figure}[t]
  \centering
  \includegraphics[width=0.84\textwidth]{ch08/fig14_sharpe_se}
  \caption{Years of independent monthly returns needed before the standard error of the Sharpe ratio,
  from \cref{eq8:sharpese}, falls to half the true ratio, against the true ratio in annual units. The
  solid line uses the market's monthly skewness and excess kurtosis, the dashed line skewness
  \prAltSkew{} and excess kurtosis \prAltKurt; the dot is the market.}
  \label{fig8:sharpese}
\end{figure}

Annualising by $\sqrt{12}$ assumes the monthly returns are uncorrelated. \textcite{lo2002} shows that
with autocorrelations $\rho_k$ the annual Sharpe ratio is $\eta\,\theta$ with
\begin{equation}\label{eq8:loeta}
  \eta=\frac{q}{\sqrt{q+2\sum_{k=1}^{q-1}(q-k)\rho_k}},\qquad q=12 .
\end{equation}
For the market $\eta=\prEta$ against $\sqrt{12}=3.46$, and the adjusted ratio is \prSRlo. The
correction is small for a traded index and large for the smoothed returns of \cref{sec8:smoothing}.

The ratios disagree when risk is not all market risk. The high technology industry portfolio of
\textcite{frenchlib} has a Sharpe ratio of \prHtSR{} against the market's \prSR, with volatility of
\prHtVol\% against \prMktVol\% and a beta of \prHtBeta. Its Treynor ratio, \prHtTreynor\%, beats the
market's \prMktTreynor\%, and its alpha is a positive \prHtAlpha\% a year. Yet the residual risk is
\prHtTE\% a year, so the information ratio is only \prHtIR, and $M^2$ is \prHtMsq\% against the
market's \prMktMean\%. An investor who holds the industry alone bears risk the market does not pay
for, and the Sharpe ratio and $M^2$ say so. For an investor adding a small position to a diversified
portfolio, beta and alpha matter more, provided the market captures the systematic risk that matters to
that investor. On the downside measures, with a zero target, the industry's annualised Sortino ratio is
\prHtSortino{} against \prMktSortino{} and its monthly Omega \prHtOmega{} against \prMktOmega.

\begin{figure}[t]
  \centering
  \includegraphics[width=0.94\textwidth]{ch08/fig06_ratios}
  \caption{Left: three laws of monthly returns with the same mean and volatility, normal and
  shifted exponentials skewed to the left and to the right. Right: the Sharpe, Sortino and Omega
  ratios and the MPPM of the two skewed laws as multiples of their values for the normal law.}
  \label{fig8:ratios}
\end{figure}

The Sharpe ratio uses two moments and cannot see the rest. \Cref{fig8:ratios} shows three laws of
monthly returns, each with a mean of \thMu\% and a volatility of \thVol\% a year and so a Sharpe ratio
of \thSR: a normal law, and exponential laws shifted and scaled to the same moments with the long tail
on the left or on the right. Their 99\% monthly VaR is \thNormalVar\%, \thLeftVar\% and \thRightVar\%,
and their ES \thNormalEs\%, \thLeftEs\% and \thRightEs\%. The Sortino ratio separates them, at
\thNormalSo, \thLeftSo{} and \thRightSo. Omega at a zero threshold barely moves, at \thNormalOm,
\thLeftOm{} and \thRightOm, because by \cref{prop8:m2}(ii) it sees only the mean and the first lower
partial moment, which are close for the three laws. The manipulation-proof measure of
\textcite{goetzmann2007},
\begin{equation}\label{eq8:mppm}
  \Theta=\frac{1}{(1-\rho)\,\Delta t}\,\ln\Bigl(\frac1T\sum_{t=1}^{T}
  \Bigl[\frac{1+r_t}{1+r_{f,t}}\Bigr]^{1-\rho}\Bigr),
\end{equation}
is an annualised certainty equivalent for an investor with constant relative risk aversion $\rho$.
With $\rho=3$ and a zero risk free rate it ranks the right skewed law first at \thRightMp\%, the normal
second at \thNormalMp\% and the left skewed last at \thLeftMp\%. The normal and the left
skewed law put a tiny probability on returns below $-100\%$, about $\thLeftTail$ for the left skewed
law. Returns at or below $-100\%$ leave the positive-wealth domain of power utility, where
$(1+r)^{-2}$ cannot be integrated, so for the untruncated normal and left skewed laws the measure is
not well defined. The finite values quoted apply only to those two laws truncated at a monthly return
of $-\thFloor\%$ and renormalised, and are computed by quadrature; the right skewed law is bounded below
and needs no cut. A sample of a million
draws gives \thRightMpSample\%, \thNormalMpSample\% and \thLeftMpSample\%. A manager who sells out of the
money puts to raise the Sharpe ratio moves the return law toward the left skewed case, and $\Theta$
falls; that is the sense in which it resists manipulation, within its assumption of a power utility
investor, and it says nothing about risks the record does not show.

\FloatBarrier
% =====================================================================================
\section{Drawdowns}
\label{sec8:drawdown}

Investors feel losses as a path, from a peak to a trough and then the long climb back. With wealth
$W_t$ and running peak $P_t=\max_{s\le t}W_s$, the \emph{drawdown} is $D_t=1-W_t/P_t$, and the maximum
drawdown over a period is the largest $D_t$ in it. The US market's maximum drawdown in daily data
from \tdFirst{} to \tdLast{} is \ddMax\%, from a peak on \ddPeak{} to a trough on \ddTrough. It took
\ddDownYears{} years to fall and \ddAllYears{} years from the peak until wealth regained it, on
\ddRecover. The next two are \ddSecond\%, from \ddSecondPeak{} to \ddSecondTrough, and \ddThird\%, from
\ddThirdPeak{} to \ddThirdTrough. The market was below a previous peak on \ddShareUnder\% of all days,
more than 10\% below on \ddShareTen\% and more than 20\% below on \ddShareTwenty\%. The compound growth
rate over the whole period is \ddCagr\% a year, and its ratio to the maximum drawdown is \ddCalmar.
This is a whole period form of the \emph{Calmar ratio}\index{Calmar ratio}, which \textcite{young1991} defined over the
last 36 months and which in practice is still usually computed that way, so it gives a different number
for the same fund.

\begin{figure}[t]
  \centering
  \includegraphics[width=0.94\textwidth]{ch08/fig07_drawdown}
  \caption{Left: drawdown of the US market from its running peak, daily. Right: the share of days
  on which the log drawdown $\ln(P_t/W_t)$ exceeded a given level, against the exponential law of
  \cref{prop8:ddlaw} with the sample's log drift and volatility, on a logarithmic scale.}
  \label{fig8:drawdown}
\end{figure}

A drawdown depends on the path and so on the order of returns, which makes it hard to work with
directly. For a random walk in logs it has a simple long run law.

\begin{proposition}[The law of the drawdown of a random walk]\label{prop8:ddlaw}
Let $X_t=\ln W_t=\mu t+\sigma B_t$ with $B$ a standard Brownian motion and $\mu>0$, and let
$Y_t=\max_{s\le t}X_s-X_t$ be the log drawdown. As $t\to\infty$ the law of $Y_t$ converges to the
exponential law with mean $\sigma^2/(2\mu)$, so that
$\mathbb P(D_t>d)\to(1-d)^{2\mu/\sigma^2}$ for $0<d<1$.
\end{proposition}
\begin{proof}
For fixed $t$, $Y_t=\max_{0\le u\le t}\bigl(X_{t-u}-X_t\bigr)=\max_{0\le u\le t}(-\tilde X_u)$ with
$\tilde X_u=X_t-X_{t-u}$. The increments of $\tilde X$ over $[0,t]$ are those of $X$ in reverse order, so
$\tilde X$ is again a Brownian motion with drift $\mu$ and volatility $\sigma$, and $Y_t$ has the law
of $\max_{u\le t}Z_u$ for $Z_u=-\mu u+\sigma B_u$. That maximum increases with $t$ to
$Z^*=\sup_{u\ge0}Z_u$, so $Y_t$ converges in law to $Z^*$. With $\lambda=2\mu/\sigma^2$ the process
$e^{\lambda Z_u}$ is a martingale, because $-\lambda\mu+\lambda^2\sigma^2/2=0$. For $a>0$ let $H_a$ be
the first time $Z$ reaches $a$. The stopped martingale $e^{\lambda Z_{u\wedge H_a}}$ is bounded by
$e^{\lambda a}$ and tends to $e^{\lambda a}$ on $\{H_a<\infty\}$ and to zero elsewhere, because $Z_u\to-\infty$.
Optional stopping gives $1=e^{\lambda a}\mathbb P(H_a<\infty)$, so $\mathbb P(Z^*\ge a)=e^{-\lambda a}$, the
exponential law with mean $1/\lambda=\sigma^2/(2\mu)$. Finally $D_t>d$ exactly when
$Y_t>-\ln(1-d)$.
\end{proof}

The result is standard in the theory of reflected Brownian motion \parencite{harrison1985}, and
\textcite{magdon2004} give the law of the maximum drawdown over a finite horizon. With the market's
daily log drift of \ddLogMu\% a year and log volatility of \ddLogVol\%, the law gives a mean log
drawdown of \ddLawMean\% and a probability of \ddLawTwenty\% of being more than 20\% below the peak.
In the data the mean is \ddEmpMean\% and the share of days more than 20\% down is \ddShareTwenty\%.
In \cref{fig8:drawdown} the data lie a little below the law up to log drawdowns of about 40\%, rise
above it between log drawdowns of 40\% and 55\%, simple drawdowns of about \ddSimpleLo\% to
\ddSimpleHi\%, which the declines of 2000 to 2002 and 2007 to 2009 passed through on their way down,
and fall away beyond, since one sample holds only a handful of deep episodes.

The maximum drawdown over a fixed horizon is what investors watch. In \ddSimM{} simulated ten year
paths of daily normal log returns with the market's drift and volatility, the median maximum
drawdown is \ddSimMedian\% and the 90th percentile \ddSimNinety\%; the maximum drawdown exceeds 30\%
on \ddSimThirty\% of paths and 50\% on \ddSimFifty\%. In the \ddWinN{} ten year windows of the actual
series that start a year apart, which overlap, the median is \ddWinMedian\%, the smallest \ddWinMin\%
and the largest \ddWinMax\%, and \ddWinAboveNinety{} of the \ddWinN{} lie above the simulated 90th
percentile (\cref{fig8:tenyear}). In this sample real drawdowns were deeper than the random walk predicts. That is consistent with the
clustered losses of \cref{sec8:garch} and the heavy tails of \cref{sec8:var}, but the comparison does
not separate the two, and overlapping windows built on a few episodes make it loose. An investor who reads a 30\% maximum drawdown in a ten year
record as evidence of a risky strategy should know that a broad stock index with no skill and no
leverage did worse than that in a typical decade of this sample.

\begin{figure}[t]
  \centering
  \includegraphics[width=0.84\textwidth]{ch08/fig15_ten_year}
  \caption{The maximum drawdown over ten years in \ddSimM{} simulated paths of daily normal log
  returns with the market's drift and volatility (histogram), and in the \ddWinN{} ten year windows of
  the actual series that start a year apart (marks below the axis). The dotted line is the simulated
  90th percentile. The windows overlap, so the marks are far from independent draws, and marks
  coincide where windows share the same episode.}
  \label{fig8:tenyear}
\end{figure}

\FloatBarrier
% =====================================================================================
\section{When returns are smoothed}
\label{sec8:smoothing}

Chapter~1 showed that appraisal based returns are smoothed, and Chapter~7 met the same effect in
illiquid assets. The model of \textcite{geltner1993} writes the reported return as
\begin{equation}\label{eq8:smooth}
  r^*_t=(1-\phi)\,r_t+\phi\,r^*_{t-1},\qquad 0\le\phi<1,
\end{equation}
a weighted average of the true return $r_t$ and last period's reported return.

\begin{proposition}[What smoothing does to measured risk]\label{prop8:smooth}
Let the true returns $r_t$ be uncorrelated with mean $\mu$ and variance $\sigma^2$, and let $r^*$ be the
stationary solution of \cref{eq8:smooth}. Then:
\begin{enumerate}[label=(\roman*)]
  \item $r^*$ has mean $\mu$ and variance $\sigma^2(1-\phi)/(1+\phi)$, so its volatility is understated
  by the factor $\sqrt{(1-\phi)/(1+\phi)}$ and its Sharpe ratio overstated by the inverse;
  \item the autocorrelation of $r^*$ at lag $k$ is $\phi^k$;
  \item the sum of all autocovariances of $r^*$ is $\sigma^2$, so the variance of a sum of $n$ reported
  returns, divided by $n$, tends to $\sigma^2$ as $n\to\infty$;
  \item the true return is recovered exactly by $r_t=(r^*_t-\phi r^*_{t-1})/(1-\phi)$.
\end{enumerate}
\end{proposition}
\begin{proof}
The stationary solution is $r^*_t=(1-\phi)\sum_{j\ge0}\phi^j r_{t-j}$, whose weights sum to one, so its
mean is $\mu$. Its autocovariance at lag $k\ge0$ is
$(1-\phi)^2\sigma^2\sum_{j\ge0}\phi^j\phi^{j+k}=\sigma^2\phi^k(1-\phi)^2/(1-\phi^2)=\sigma^2\phi^k(1-\phi)/(1+\phi)$,
which gives (i) at $k=0$ and (ii) on dividing. For (iii), the autocovariances at lags
$k=-\infty,\dots,\infty$ sum to $\sigma^2\frac{1-\phi}{1+\phi}\cdot\frac{1+\phi}{1-\phi}=\sigma^2$, and the
variance of a sum of $n$ terms divided by $n$ tends to that sum. (iv) is \cref{eq8:smooth} solved for
$r_t$.
\end{proof}

Smoothing hides risk at short horizons and none of it at long ones, which is why the variance ratio
of Chapter~1 detects it. To see the size of the effect, smooth the monthly total returns of the US
market with $\phi=\smPhi$, for which (i) gives a volatility ratio of \smVolRatioTheory. The measured
volatility falls from \smVolT\% to \smVolS\%, almost exactly half, and the Sharpe ratio rises from \smSRT{} to \smSRS. It is computed on returns less
the unsmoothed bill rate, so it divides mean excess returns of \smExMeanT\% and \smExMeanS\% a year by
the volatilities of those excess returns, \smExVolT\% and \smExVolS\%, which differ a little from the
volatilities of total returns. The first order autocorrelation, \smAcT{} in the market, is \smAcS{} in the
smoothed series. The 99\% monthly VaR falls from \smVarT\% to \smVarS\%. The maximum drawdown in
monthly data falls less, from \smMddT\% to \smMddS\%, because a drawdown is a sum of many returns and
by (iii) smoothing does not reduce long horizon variance. For the same reason the volatility of
calendar year returns falls only from \smAnnVolT\% to \smAnnVolS\%.

\begin{figure}[t]
  \centering
  \includegraphics[width=0.94\textwidth]{ch08/fig08_smoothing}
  \caption{Left: monthly returns of the US market and the same returns smoothed with $\phi=\smPhi$,
  2007 to 2010. Right: volatility a year, 99\% monthly VaR and maximum drawdown of the market, the
  smoothed series and the series unsmoothed with the estimated $\hat\phi$.}
  \label{fig8:smoothing}
\end{figure}

Three corrections are available. Lo's factor \cref{eq8:loeta} applied to the smoothed series gives a
Sharpe ratio of \smSRloS, close to the \smSRloT{} that the same correction gives for the true series,
because twelve lags capture most of the autocorrelation. Regressing the smoothed returns on the market
gives a beta of \smBetaZero, close to $1-\phi$; adding three lags of the market and summing the
coefficients, as \textcite{dimson1979} proposed for thinly traded shares, gives \smBetaSum, against
$1-\phi^4=0.87$ in theory. Inverting with
\cref{prop8:smooth}(iv) and the estimate $\hat\phi=\smPhiHat$ gives back a series whose correlation
with the true returns is \smCorrU, with volatility \smVolU\%, Sharpe ratio \smSRU, VaR \smVarU\% and
maximum drawdown \smMddU\% (\cref{fig8:smoothing}). Real appraisal processes are not exactly
\cref{eq8:smooth}, and $\phi$ must be estimated from a short record, so the inversion is a correction
with error of its own, and an inversion based on the wrong model of smoothing need not improve the
figures. A risk figure computed from smoothed returns should therefore come with a test for
autocorrelation and, where the model is credible, a corrected value beside it. The R listing below
smooths the monthly market returns with $\phi=\smPhi$, estimates $\phi$ by the lag one autocorrelation,
inverts, and fits the contemporaneous and Dimson betas. It prints $\hat\phi$, the correlation of the
inverted series with the true one, the three volatilities and the two betas quoted in this section.
\Needspace*{16\baselineskip}
\lstinputlisting[language=R]{code/ch08/snippets/unsmooth.R}

\FloatBarrier
% =====================================================================================
\section{Reporting performance to clients}
\label{sec8:gips}

A manager's track record is usually a \emph{composite}: the combined return of all the portfolios
it manages to one strategy. The Global Investment Performance Standards \parencite{gips2019} set out
how firms that claim compliance build composites and report them. All actual, fee paying,
discretionary portfolios must belong to at least one composite, so a firm cannot show only its best
accounts, and portfolios that leave a composite stay in its history up to the date they left.
Composite returns are asset weighted, either with beginning of period values or with values that
also reflect external cash flows, or are computed by the aggregate method, which treats the composite
as one portfolio.

\begin{proposition}[Two ways to compute a composite return agree]\label{prop8:composite}
Let portfolio $i$ have value $V_i$ at the start of a period with no external cash flows and return
$r_i$ over it. The return weighted by beginning values, $\sum_iw_ir_i$ with $w_i=V_i/\sum_jV_j$, equals the
return of the aggregate, $\bigl(\sum_iV_i(1+r_i)-\sum_iV_i\bigr)/\sum_iV_i$.
\end{proposition}
\begin{proof}
The numerator of the aggregate return is $\sum_iV_ir_i$, and dividing by $\sum_jV_j$ gives
$\sum_iw_ir_i$.
\end{proof}

With external cash flows, simple beginning of period weights no longer reproduce the aggregate return,
and each method has to treat the timing of the flows consistently; within each portfolio the time
weighted methods of Chapter~1 apply. Beyond returns, the standards ask for a measure of \emph{internal dispersion}, how far
the returns of portfolios within the composite spread around it, such as their standard deviation over
the year; it is not required when the composite holds five or fewer portfolios for the full year. Once 36
monthly returns are available, they also ask for the three year annualised ex post standard deviation
of the composite and its benchmark, computed from those monthly returns.

Dispersion matters because an equal weighted and an asset weighted view of the same composite can
disagree. If a large account did badly and two small ones did well, the asset weighted return is low
and the equal weighted return high; \cref{exr8:eight} works through a case. A prospective client
should ask which view the headline figure takes and how its own account would have fared.

The tools of this chapter give a short list of questions for reviewing a manager. How long is the
record, and what is the standard error of its Sharpe ratio (\cref{prop8:sharpese})? Are the returns
autocorrelated, and if so what do the corrections of \cref{sec8:smoothing} do to the volatility? What
was the maximum drawdown, and how does it compare with that of a passive index over the same years? Is
the return distribution skewed to the left, as it is for strategies that sell insurance, and what do
the Sortino ratio and the MPPM say? Does the composite include every account managed to the strategy,
and how widely did those accounts differ?

\FloatBarrier
% =====================================================================================
\section{Case study: seven risk models on two markets}
\label{sec8:case}

This section runs seven daily VaR models through the tests of \cref{sec8:backtest} on the US market
and on the Canadian dollar. Each model forecasts the 99\% VaR and the 97.5\% ES of the next day's loss
using only returns up to the close of the day before:
\begin{itemize}
  \item \emph{normal}: the mean and standard deviation of the last 250 days in the normal formulas;
  \item \emph{historical}: the empirical quantile and tail mean of the last 250 days, and again of the
        last 1,000;
  \item \emph{EWMA}: the RiskMetrics variance with $\lambda=0.94$ and the normal law;
  \item \emph{GARCH, normal} and \emph{GARCH, $t$}: \cref{eq8:garch} with normal or Student $t$
        innovations, refitted every 250 days on all the data to date;
  \item \emph{filtered historical}: the normal GARCH volatility forecast times the empirical quantile and tail
        mean of its own standardised residuals to date \parencite{barone1999,hull1998}.
\end{itemize}
The first 1,000 days are the initial estimation sample and are not scored. Every refit and every
residual quantile uses only returns to the close of the day before the forecast, and between refits the
GARCH parameters are held fixed while the variance is updated each day; the historical windows roll
forward daily. For the US market the test period runs from
\caFirst{} to \tdLast, \caN{} days with \caExpected{} exceptions expected at 99\%. For a Canadian
investor holding US dollars, the loss is minus the daily log change of the USD/CAD rate, quoted in
Canadian dollars per US dollar, and the
test period runs from \cyFirst{} to \cyLast, \cyN{} days with \cyExpected{} expected.

\begin{table}[t]
  \centering
  \small
  \caption{Backtests of seven daily VaR models. Exceptions are days on which the loss exceeded the 99\%
  VaR; $p_K$ and $p_C$ are the $p$ values of the Kupiec test and the Christoffersen\index{Christoffersen} independence
  test; red is the share of 250 day windows with ten or more exceptions, in percent; $Z_2$ is
  \cref{eq8:z2} at 97.5\%, whose approximate 5\% critical value under each model's own law, treating the fitted parameters as
  fixed, is between \caCritNor{} and \caCritEmp{} (see the text). These are reference values and do
  not allow for the estimation in the rolling procedure.}
  \label{tab8:case}
  \setlength{\tabcolsep}{4pt}
  \begin{tabular}{@{}lrrrrrrrr@{}}
    \toprule
    & \multicolumn{5}{c}{US market, \caExpected{} expected} & \multicolumn{3}{c}{USD/CAD, \cyExpected{} expected}\\
    \cmidrule(lr){2-6}\cmidrule(l){7-9}
    Model & Exceptions & $p_K$ & $p_C$ & Red & $Z_2$ & Exceptions & $p_K$ & $p_C$\\
    \midrule
    \caTableRows
    \bottomrule
  \end{tabular}
\end{table}

\Cref{tab8:case} gives the results. On the US market every model has too many exceptions, and with
\caN{} days the Kupiec test rejects them all. The normal model on a 250 day window is worst, at
\caNorX{} exceptions, a rate of \caNorRate\%, with \caNorRed\% of its 250 day windows in the red zone and as many
as \caMaxNor{} exceptions in one 250 day window. Historical simulation does better on the count, at
\caHsX{} for the short window and \caHslX{} for the long one, but its exceptions bunch: it reaches
\caMaxHs{} in one window, and the Christoffersen test rejects independence at any usual level. A window
of past returns adapts to a new level of volatility only as the old days drop out of it.

The GARCH models reduce the bunching without removing it: the Christoffersen test still rejects
independence for both, with $p_C$ of \caGnPc{} for normal GARCH and \caGtPc{} for the $t$ version.
Normal GARCH still has \caGnX{} exceptions, because the normal quantile is too thin for residuals whose
1\% quantile is $-\gaResQ$. The $t$ version brings the count to \caGtX, with no window in the red. On
the US market filtered historical simulation is best on every measure, though it too fails the count
test: \caFhsX{} exceptions, a rate of \caFhsRate\% against 1\%, no red windows, at most
\caMaxFhs{} exceptions in any 250 days, a Christoffersen $p$ value of \caFhsPc, the only one above 5\%,
and the smallest $Z_2$. It combines the two ideas of this chapter, a variance that responds to clustering and a tail
taken from the data. \Cref{fig8:case} shows the difference over time. The crash of 19 October 1987 was
\caOctHs{} times the historical VaR of that morning and \caOctFhs{} times the filtered historical VaR;
no daily model forecasts a day like that.

\Cref{fig8:crisis} shows why the window model bunches its exceptions. Its VaR moves only when a large
loss enters or leaves the last 250 days. Over 2007 to 2009 it reached its peak of \crHsMax\% only on
\crHsPeak, and it held that level until \crHsDrop, the day on which the loss of \crDropDate, equal to
the peak, left the window. The filtered VaR multiplies a quantile of past
standardised residuals by the current volatility forecast, and it peaked on \crFhsPeak. In the
\crDays{} trading days of the three years, with \crExpected{} exceptions expected, historical simulation had \crHs{} and filtered
historical simulation \crFhs; in the \crAutDays{} days from September to November 2008 they had
\crAutHs{} and \crAutFhs.

\begin{figure}[t]
  \centering
  \includegraphics[width=0.94\textwidth]{ch08/fig09_case}
  \caption{Exceptions of the 99\% VaR in the previous 250 trading days for historical simulation on a
  250 day window and filtered historical simulation, US market. The shaded bands are the green
  zone, four exceptions or fewer, and the red zone, ten or more.}
  \label{fig8:case}
\end{figure}

\begin{figure}[t]
  \centering
  \includegraphics[width=0.94\textwidth]{ch08/fig16_crisis}
  \caption{Daily losses of the US market, 2007 to 2009 (grey), with the 99\% VaR of historical
  simulation on 250 days and of filtered historical simulation, as scored in \cref{tab8:case}. The
  marks below zero are the exceptions of each model, historical above and filtered below.}
  \label{fig8:crisis}
\end{figure}

The ES statistic points the same way. Each model's $Z_2$ is compared with its distribution under
the null, worked out by the calculation of \cref{sec8:backtest} for the loss law the model assumes over
the \caN{} days of the test. With normal innovations, for the normal, EWMA and normal GARCH models, the 5\% critical value is
\caCritNor; with $t$ innovations and the \caNuLast{} degrees of freedom of the last refit it is
\caCritT; and under the empirical law of the GARCH residuals that filtered historical simulation uses it
is \caCritEmp, which also serves as a guide for plain historical simulation, whose null has no fixed
law. The smallest $Z_2$ in \cref{tab8:case}, \caFhsEs, is about three null standard deviations above zero
and about twice the largest of these reference values. That indicates that all seven models
underpredict the 97.5\% shortfall of the US market relative to these fixed-parameter benchmarks, but
it does not establish rejection at a calibrated significance level for the rolling forecasting
procedures. A calibrated test would draw histories from each fitted model and repeat the
estimation, the refits and the volatility updates on every one; that simulation is not run here. At 97.5\% the exceptions, against
\caExpectedEs{} expected, are \caNorXs{} for the normal model, \caHsXs{} for historical simulation and
\caFhsXs{} for filtered historical simulation; part of each $Z_2$ comes from those excess exceptions,
and the rest from losses that are larger than forecast when they occur.

On the USD/CAD sample the ranking changes. Its exceptions bunch less, and no model fails the
Christoffersen test at 5\%. Historical simulation on 250 days, with \cyHsX{} exceptions against
\cyExpected, and normal GARCH, with \cyGnX, pass the Kupiec test. The $t$ GARCH model, the second best on
the US market, is now too cautious, with \cyGtX{} exceptions and a Kupiec $p$ value below 0.001, and
filtered historical simulation, with \cyFhsX, is rejected at 5\% for the same reason. These results hold for
these samples and these loss conventions. A model that suits one series can be wrong for another, so a
backtest has to be run on the portfolio the model will be used for.

\FloatBarrier
% =====================================================================================
\section{Limitations}
\label{sec8:limits}

Every result in this chapter rests on a model, and each model leaves something out.
\begin{itemize}
  \item VaR says nothing about losses beyond it, and ES depends on a tail that a few observations
        determine. Both are statistics of one day's loss on a fixed position; liquidating a position
        takes time, and a ten day horizon does not capture a market that stays closed or one-sided.
  \item The square root of time rule and \cref{prop8:garch}(ii) assume a stable model. After the
        crises in the sample, volatility fell faster than the fitted persistence implied.
  \item Extreme value estimates rely on a threshold choice and on the assumption that the tail of
        the past is the tail of the future. The simulation of \cref{sec8:evt} shows the shape is
        biased at a 95\% threshold even when the law is known.
  \item The backtests assume a fixed portfolio. A trading book changes daily, and its exceptions
        mix model error with the effect of trades made during the day.
  \item The performance ratios assume that returns are measured fairly. Smoothing, stale prices and
        selective composites bias them all, and a short record makes any ratio imprecise.
  \item The drawdown law of \cref{prop8:ddlaw} assumes independent normal increments, and the
        data show deeper declines.
  \item The data are one equity index and one exchange rate. Credit, options and illiquid assets
        have risks these models were not built for.
\end{itemize}

\FloatBarrier
% =====================================================================================
\section{Summary}
\label{sec8:summary}

\begin{itemize}
  \item VaR is a quantile and can fail subadditivity; ES is the average loss beyond it, is coherent
        and has a representation as a worst case over scenarios. For daily equity returns the normal
        law understates both, and Cornish-Fisher can overcorrect.
  \item GARCH(1,1) captures volatility clustering. Its forecasts revert to the long run variance
        at the rate $a+b$, and with normal innovations its kurtosis exceeds three.
  \item The generalised Pareto tail gives VaR and ES in closed form beyond the range of the data, and in
        the simulation here it estimated a high quantile with less spread than the sample did, even with
        a biased shape.
  \item The hits of a correct VaR model are independent Bernoulli variables. Kupiec tests the count
        and Christoffersen the independence; both have little power over one year.
  \item The Sharpe ratio has a standard error that grows with negative skew and fat tails, and its
        annual value needs Lo's correction when returns are autocorrelated. Treynor, alpha, $M^2$,
        Sortino, Omega and the MPPM answer different questions.
  \item The log drawdown of a random walk with drift has an exponential long run law with mean
        $\sigma^2/(2\mu)$.
  \item Smoothing divides measured volatility by $\sqrt{(1+\phi)/(1-\phi)}$ and leaves long run variance
        unchanged, and the Geltner inversion recovers the true series exactly when $\phi$ is known and
        closely when it is estimated.
  \item On the US market from 1967, filtered historical simulation was the only model to pass the
        independence test; all seven failed the Kupiec count test, and all had $Z_2$ statistics above the
        fixed-parameter reference values, a comparison that is not a fully calibrated test of the
        rolling procedures.
\end{itemize}

\FloatBarrier
% =====================================================================================
\section{Exercises}
\label{sec8:exercises}

\begin{enumerate}[label=\textbf{8.\arabic*}, leftmargin=*, ref=8.\arabic*]
  \item\label{exr8:one} A position of 10 million has normal daily returns with mean zero and
  volatility 1.2\%. Find its one day 99\% VaR, its one day 97.5\% ES and its ten day 99\% VaR by the
  square root of time rule.
  \item\label{exr8:two} Two independent bonds each lose 100 on default, with probability 4\%. At
  95\%, find the VaR of one bond and of the pair, and show that VaR is not subadditive here. Then find
  the ES of one bond and of the pair and check subadditivity.
  \item\label{exr8:three} A GARCH(1,1) model in daily percent has $\omega=0.02$, $a=0.08$ and $b=0.90$,
  and tomorrow's variance is 4. Find the long run variance and annual volatility, the variance forecast
  ten days ahead, the half life of a shock and the ten day volatility, and compare the last with the
  square root of time rule.
  \item\label{exr8:four} A 99\% VaR model has 8 exceptions in 250 days. Find the Kupiec statistic and
  its $p$ value, the zone, and the probability of 8 or more exceptions under a correct model.
  \item\label{exr8:five} Of $n=10{,}000$ daily losses, $n_u=500$ exceed $u=2\%$, and the excesses fit
  a generalised Pareto law with $\xi=0.25$ and $\beta=0.7\%$. Find the 99.9\% VaR and ES.
  \item\label{exr8:six} Over $T=120$ months a fund has a monthly Sharpe ratio of 0.15, skewness $-1$ and
  excess kurtosis 3. Find the standard error of the monthly Sharpe ratio under normality and in
  general, the annualised ratio and its standard error, and the $t$ statistic.
  \item\label{exr8:seven} A fund's true returns are uncorrelated with volatility 16\% a year, and it
  reports returns smoothed with $\phi=0.5$. Find the measured volatility and the factor by which the
  Sharpe ratio is overstated.
  \item\label{exr8:eight} A composite holds three accounts worth 10, 30 and 60 at the start of the year,
  returning 5\%, 2\% and $-1\%$. Find the asset weighted and equal weighted composite returns and the
  equal weighted and asset weighted standard deviations of the account returns.
  \item\label{exr8:nine} A portfolio's log value has drift 6\% and volatility 15\% a year. Find the mean
  of its long run log drawdown and the probability that it is more than 30\% below its peak.
  \item\label{exr8:ten} A fund has an excess return of 8\% a year, volatility 20\% and beta 1.1, and the
  market has an excess return of 6\% and volatility 15\%. Find the fund's $M^2$ and Treynor ratio and
  compare them with the market's.
\end{enumerate}

\subsection*{Computing projects}
\label{sec8:projects}
These use the data and code in the repository. Each is solved by a function in
\texttt{code/ch08/projects.py}, which \texttt{test\_projects.py} runs, so an answer can be
checked with \texttt{python code/ch08/projects.py} from the root of the repository. The
shipped series end in September 2026, and those that end earlier are noted.

\begin{enumerate}[label=\textbf{P8.\arabic*}, leftmargin=*, ref=P8.\arabic*]
\item Forecast the one day 99\% VaR of a long US dollar position against the Canadian dollar from 2012 with
normal and Student t GARCH(1,1) refitted each year on the previous 1,000 days, with EWMA and with
historical simulation. Count exceptions and apply \texttt{kupiec} and \texttt{christoffersen}. Which models
pass, and which fail by too many or too few exceptions?
\item For each of the six markets of \texttt{market\_returns}, report the annualised Sharpe ratio, its standard
error from \texttt{sharpe\_se} with and without skewness and excess kurtosis, and the largest drawdown.
Which Sharpe ratios are more than two standard errors from zero, and for which market does the moment
correction matter most?
\end{enumerate}

\subsection*{Answers}

\begin{enumerate}[label=\textbf{8.\arabic*}, leftmargin=*]
  \item The VaR is $2.326\times0.012\times10=\xaVar$ million and the ES
  $\varphi(1.960)/0.025\times0.012\times10=\xaEs$ million, close to the VaR, as the ratio of
  \esRatioNormal{} in \cref{sec8:measures} leads one to expect. The ten day VaR is $\sqrt{10}\times\xaVar=\xaVarTen$ million.
  \item The chance of a default is 4\%, below 5\%, so the VaR of one bond is zero. The chance that at
  least one of the pair defaults is \xbJoint\%, so the VaR of the pair is 100, more than the sum of the
  two VaRs. The ES of one bond is $0.04\times100/0.05=\xbEsOne$. For the pair, both default with
  probability \xbBoth\%, so the worst 5\% of outcomes are a loss of 200 with probability 0.16\% and of 100
  with probability 4.84\%, and the ES is $(0.32+4.84)/0.05=\xbEsTwo$, less than $2\times\xbEsOne$.
  \item The long run variance is $0.02/0.02=\xcLr$, an annual volatility of $\sqrt{252}=\xcLrVol\%$. By
  \cref{prop8:garch}(ii) the forecast ten days ahead is $1+0.98^9\times3=\xcTen$. The half life is
  $\ln0.5/\ln0.98=\xcHalf$ days. If the daily returns are uncorrelated, the variance of the ten day return is the sum of the ten
  forecasts, and its square
  root is \xcTenDay\%, against $\sqrt{40}=\xcTenSqrt\%$ by the square root rule, which ignores the
  reversion toward the long run.
  \item With $\hat\pi=0.032$, \cref{eq8:kupiec} gives $\mathrm{LR}_{\mathrm{uc}}=\xdLR$ and $p=\xdP$, so
  the test rejects at 1\%. Eight exceptions are in the yellow zone. Under a correct model eight or more
  occur with probability \xdTail\%.
  \item By \cref{prop8:gpd}, $\mathrm{VaR}=0.02+\frac{0.007}{0.25}\bigl[(20\times0.001)^{-0.25}-1\bigr]=\xeVar\%$,
  and $\mathrm{ES}=(\xeVar\%+0.7\%-0.5\%)/0.75=\xeEs\%$.
  \item By \cref{eq8:sharpese} the standard error is $\sqrt{1.01125/120}=\xfIid$ under normality and
  $\sqrt{(1.01125+0.15+0.016875)/120}=\xfNn$ in general. The annualised ratio is $\sqrt{12}\times0.15=\xfAnn$
  with standard error \xfAnnSe, and the $t$ statistic is \xfT, short of two even after ten years.
  \item By \cref{prop8:smooth}(i) the measured volatility is $16\sqrt{0.5/1.5}=\xgVol\%$, and the Sharpe
  ratio is overstated by $\sqrt3=\xgInfl$.
  \item The asset weighted return is $(0.5+0.6-0.6)/100=\xhAw\%$ and the equal weighted return
  \xhEw\%. The equal weighted standard deviation, with divisor three, is \xhDisp{} points and the
  asset weighted one \xhDispAw{} points. The large account's loss pulls the asset weighted return
  down.
  \item By \cref{prop8:ddlaw} the mean is $0.15^2/(2\times0.06)=\xiMean\%$ in logs. A drawdown of more
  than 30\% is a log drawdown of more than $-\ln0.7=\xiLogThirty$, with probability
  $e^{-\xiLogThirty/0.1875}=\xiThirty\%$.
  \item $M^2=8\times15/20=\xjMsq\%$, equal to the market's 6\%, so the two have the same Sharpe ratio of
  0.40. The Treynor ratio is $8/1.1=\xjTreynor\%$ against the market's 6\%, so per unit of beta the fund
  did better, and the extra volatility was risk that diversification would remove.
\end{enumerate}

\FloatBarrier
% =====================================================================================
\section{Notes, sources and credits}
\label{sec8:credits}

\paragraph{Sources}
Coherent risk measures are from \textcite{artzner1999} and the coherence of expected shortfall from
\textcite{acerbi2002}. ES is shown not to be elicitable by \textcite{gneiting2011}, and its
backtests are discussed by \textcite{acerbi2014}. The Cornish-Fisher expansion is from \textcite{cornish1938}. ARCH is from
\textcite{engle1982}, GARCH from \textcite{bollerslev1986} and the $t$ innovations from
\textcite{bollerslev1987}; the EWMA rule is that of \textcite{riskmetrics1996}. Filtered historical
simulation follows \textcite{barone1999} and \textcite{hull1998}. The limit theorem for exceedances is
from \textcite{balkema1974} and \textcite{pickands1975}, the Hill estimator from \textcite{hill1975},
the theory from \textcite{embrechts1997} and the GARCH and extreme value combination from
\textcite{mcneil2000}. The coverage tests are from \textcite{kupiec1995} and
\textcite{christoffersen1998}, the zones from \textcite{bcbs1996} and the 97.5\% ES standard from
\textcite{bcbs2019}. The performance ratios are from \textcite{sharpe1966,sharpe1994},
\textcite{treynor1965}, \textcite{jensen1968}, \textcite{modigliani1997}, \textcite{sortino1994},
\textcite{keating2002} and \textcite{goetzmann2007}; the Sharpe ratio's standard error and the
annualising factor are from \textcite{lo2002}, with the general case in \textcite{opdyke2007}. Joint
elicitability of VaR and ES is from \textcite{fissler2016}, and summed lagged betas from
\textcite{dimson1979}. The
drawdown law is standard for reflected Brownian motion \parencite{harrison1985}, and the finite
horizon law is in \textcite{magdon2004}. Smoothing follows \textcite{geltner1993} and
\textcite{getmansky2004}, and the stylised facts of daily returns are surveyed by \textcite{cont2001}.
The composite rules are from the GIPS standards \parencite{gips2019}.

\paragraph{Software used}
The scripts call NumPy \parencite{numpy} for arrays and random numbers, SciPy \parencite{scipy}
for distributions, optimisation and the generalised Pareto fit, pandas \parencite{pandas} to read the
data files, and Matplotlib \parencite{matplotlib} for the figures. The tests compare the GARCH fits
with the \texttt{arch} package \parencite{arch} and the autocorrelations with statsmodels
\parencite{statsmodels}. The cross-checks run in R \parencite{rcore} and GNU Octave \parencite{octave}. \texttt{code/requirements-lock.txt} records the full environment.

\paragraph{Projects for comparison}
The QuantStats package \parencite{quantstats} computes most of the performance ratios and drawdown statistics of
\cref{sec8:ratios,sec8:drawdown} from a return series, and \texttt{arch} has GARCH models with many
innovation laws and a bootstrap; both are useful for checking the code here. The full list, with
links, is in \texttt{CREDITS.md}.

\paragraph{Data}
The daily and monthly market returns are from the data library of \textcite{frenchlib}, in the release
built from the CRSP database of August 2026 and retrieved on 19 October 2026; the daily file runs from
\tdFirst{} to \tdLast. The library revises its files when the underlying database is updated, so the numbers here
belong to that release. The files are not redistributed with the code:
\texttt{scripts/build\_french.py -{}-download} saves them and builds the files the code reads. The
USD/CAD rate is the Bank of Canada series used in Chapter~5 \parencite{bocfx}, in Canadian dollars per US
dollar: the noon rate, \texttt{IEXE0101}, to 30 December 2016, joined to the daily average rate,
\texttt{FXUSDCAD}, from 3 January 2017, ending on \cyLast. The two loan example, the laws of \cref{fig8:ratios}, the simulated paths and the exercise
inputs are illustrations.

\end{chapterrefs}
\flushbottom
